% Источники: 04_Системные_требования/DOC-05_Системные_требования.md, 04_Системные_требования/Сессии_задач.md, 04_Системные_требования/Интеграция.md, 04_Системные_требования/Ранцевые_задачи.md
% Генератор: docs/build_pdf.py
% Шаблон печатного представления. Содержимое подставляет docs/build_pdf.py.
\documentclass[12pt,a4paper]{article}
\usepackage{iftex}
\ifPDFTeX
  \usepackage[utf8]{inputenc}
  \usepackage[T2A]{fontenc}
  \usepackage{newtxtext,newtxmath}
\else
  \usepackage{fontspec}
  \IfFontExistsTF{Times New Roman}
    {\setmainfont{Times New Roman}}
    {\setmainfont{Libertinus Serif}}
  \IfFontExistsTF{Arial}{\setsansfont{Arial}}{\setsansfont{Latin Modern Sans}}
  \IfFontExistsTF{Menlo}{\setmonofont{Menlo}[Scale=MatchLowercase]}
    {\setmonofont{Latin Modern Mono}[Scale=MatchLowercase]}
\fi
\usepackage[russian]{babel}
\usepackage{geometry}
\geometry{left=25mm,right=15mm,top=20mm,bottom=20mm}
\usepackage{setspace}
\onehalfspacing
\usepackage{indentfirst}
\setlength{\parindent}{1.25cm}
\setlength{\parskip}{2pt}
\setlength{\emergencystretch}{3em}
\usepackage{graphicx,booktabs,longtable,array,tabularx,calc,etoolbox}
\usepackage{enumitem}
\setlist[itemize]{leftmargin=1.25cm,itemsep=2pt,topsep=4pt}
\setlist[enumerate]{leftmargin=1.25cm,itemsep=2pt,topsep=4pt}
\newcommand{\tightlist}{\setlength{\itemsep}{0pt}\setlength{\parskip}{0pt}}
\AtBeginEnvironment{longtable}{\small\singlespacing}
\setlength{\LTpre}{8pt}
\setlength{\LTpost}{8pt}
\setlength{\tabcolsep}{5pt}
\setlength{\LTleft}{0pt}
\setlength{\LTright}{0pt}
\usepackage{xurl}
\usepackage[unicode,colorlinks=true,linkcolor=black,urlcolor=blue]{hyperref}
\usepackage{bookmark}
\hypersetup{pdftitle={Системные требования. Сессии, интеграция и ранцевые задачи},pdfauthor={Анохин Г.А.},pdfsubject={Криптотренажёр mephi-12}}
\urlstyle{same}
\pagestyle{plain}
\setcounter{secnumdepth}{3}
\setcounter{tocdepth}{2}

\begin{document}
\hypersetup{pageanchor=false}
\begin{titlepage}
  \thispagestyle{empty}
  \centering
  {\large НИЯУ МИФИ\par}
  \vspace*{3.5cm}
  {\LARGE\bfseries Системные требования. Сессии, интеграция и ранцевые задачи\par}
  \vspace{0.5cm}
  {\large DOC-05\par}
  \vspace{0.6cm}
  {\large Платформа для изучения и практики курса по криптографии\par}
  {\large Криптотренажёр mephi-12\par}
  \vspace{0.7cm}
  {\large Версия 0.4\par}
  {\large Редакция от 07.11.2026\par}
  \vspace{1cm}
  \begin{tabularx}{\textwidth}{@{}p{3.4cm}Xp{4cm}@{}}
    \toprule
    Роль & ФИО & Подпись / Дата \\
    \midrule
    Разработал & Анохин Г.А. & \\
    \addlinespace[6pt]
    Согласовано & & \\
    \addlinespace[6pt]
    Утверждено & & \\
    \bottomrule
  \end{tabularx}
  \vfill
  {\large Москва, 2026\par}
\end{titlepage}
\hypersetup{pageanchor=true}

\section*{Лист регистрации изменений}
\addcontentsline{toc}{section}{Лист регистрации изменений}
\begingroup\small\singlespacing
\noindent
\begin{tabularx}{\textwidth}{@{}p{1.2cm}p{2cm}p{2.5cm}Xp{1.8cm}@{}}
\toprule
Версия & Дата & Автор & Содержание изменения & Утвердил \\
\midrule
0.4 & 07.11.2026 & Анохин Г.А. & Уточнены правила и контракты; расширено описание сценариев; подготовлена печатная редакция. & \\
\bottomrule
\end{tabularx}
\endgroup
\clearpage
\begingroup\small\singlespacing
\tableofcontents
\endgroup
\clearpage

\hypertarget{doc-05-системные-требования}{}

{\def\LTcaptype{none} % do not increment counter
\begin{longtable}[]{@{}
  >{\raggedright\arraybackslash}p{(\linewidth - 2\tabcolsep) * \real{0.3200}}
  >{\raggedright\arraybackslash}p{(\linewidth - 2\tabcolsep) * \real{0.6800}}@{}}
\toprule\noalign{}
\begin{minipage}[b]{\linewidth}\raggedright
Поле
\end{minipage} & \begin{minipage}[b]{\linewidth}\raggedright
Значение
\end{minipage} \\
\midrule\noalign{}
\endhead
\bottomrule\noalign{}
\endlastfoot
Версия & 0.4 \\
Дата подготовки & 07.11.2026 \\
Редактор & Анохин Г.А. \\
Область редакции & Сессии задач и общесистемная интеграция \\
\end{longtable}
}

\section{Назначение и уровень описания}\label{1-ux43dux430ux437ux43dux430ux447ux435ux43dux438ux435-ux438-ux443ux440ux43eux432ux435ux43dux44c-ux43eux43fux438ux441ux430ux43dux438ux44f}

\hypertarget{1-назначение-и-уровень-описания}{}

Документ задаёт наблюдаемое поведение учебного криптотренажёра mephi-12 и свойства взаимодействия его частей. Он раскрывает \href{../02_\%D0\%A2\%D0\%B5\%D1\%85\%D0\%BD\%D0\%B8\%D1\%87\%D0\%B5\%D1\%81\%D0\%BA\%D0\%BE\%D0\%B5_\%D0\%B7\%D0\%B0\%D0\%B4\%D0\%B0\%D0\%BD\%D0\%B8\%D0\%B5/\%D0\%A2\%D0\%B5\%D1\%85\%D0\%BD\%D0\%B8\%D1\%87\%D0\%B5\%D1\%81\%D0\%BA\%D0\%BE\%D0\%B5_\%D0\%B7\%D0\%B0\%D0\%B4\%D0\%B0\%D0\%BD\%D0\%B8\%D0\%B5.pdf}{ТЗ DOC-02}. Программные обязанности определяются в \href{../07_\%D0\%9F\%D1\%80\%D0\%BE\%D0\%B3\%D1\%80\%D0\%B0\%D0\%BC\%D0\%BC\%D0\%BD\%D1\%8B\%D0\%B5_\%D1\%82\%D1\%80\%D0\%B5\%D0\%B1\%D0\%BE\%D0\%B2\%D0\%B0\%D0\%BD\%D0\%B8\%D1\%8F/DOC-06_\%D0\%9F\%D1\%80\%D0\%BE\%D0\%B3\%D1\%80\%D0\%B0\%D0\%BC\%D0\%BC\%D0\%BD\%D1\%8B\%D0\%B5_\%D1\%82\%D1\%80\%D0\%B5\%D0\%B1\%D0\%BE\%D0\%B2\%D0\%B0\%D0\%BD\%D0\%B8\%D1\%8F.pdf\#nameddest=doc-06-\%D1\%81\%D0\%B2\%D0\%BE\%D0\%B4\%D0\%BD\%D1\%8B\%D0\%B5-\%D0\%BF\%D1\%80\%D0\%BE\%D0\%B3\%D1\%80\%D0\%B0\%D0\%BC\%D0\%BC\%D0\%BD\%D1\%8B\%D0\%B5-\%D1\%82\%D1\%80\%D0\%B5\%D0\%B1\%D0\%BE\%D0\%B2\%D0\%B0\%D0\%BD\%D0\%B8\%D1\%8F}{DOC-06}; точные форматы операций сессий вынесены в \href{../07_\%D0\%9F\%D1\%80\%D0\%BE\%D0\%B3\%D1\%80\%D0\%B0\%D0\%BC\%D0\%BC\%D0\%BD\%D1\%8B\%D0\%B5_\%D1\%82\%D1\%80\%D0\%B5\%D0\%B1\%D0\%BE\%D0\%B2\%D0\%B0\%D0\%BD\%D0\%B8\%D1\%8F/DOC-06.2_\%D0\%A1\%D0\%B5\%D1\%81\%D1\%81\%D0\%B8\%D0\%B8_\%D0\%B7\%D0\%B0\%D0\%B4\%D0\%B0\%D1\%87_\%D0\%B8_\%D0\%BF\%D0\%BB\%D0\%B0\%D0\%BD\%D0\%B8\%D1\%80\%D0\%BE\%D0\%B2\%D1\%89\%D0\%B8\%D0\%BA.pdf\#nameddest=doc-062-\%D0\%BA\%D0\%BE\%D0\%BD\%D1\%82\%D1\%80\%D0\%B0\%D0\%BA\%D1\%82\%D1\%8B-\%D1\%81\%D0\%B5\%D1\%81\%D1\%81\%D0\%B8\%D0\%B9-\%D0\%B7\%D0\%B0\%D0\%B4\%D0\%B0\%D1\%87}{приложение к DOC-06.2}.

Авиационная основа комплекта --- DO-178C/КТ-178С и ARP4754. Эта редакция не содержит утверждения о полном соответствии стандартам или сертификации. Новые нормативные положения и номера пунктов здесь не устанавливаются.

\section{Система и взаимодействие частей}\label{2-ux441ux438ux441ux442ux435ux43cux430-ux438-ux432ux437ux430ux438ux43cux43eux434ux435ux439ux441ux442ux432ux438ux435-ux447ux430ux441ux442ux435ux439}

\hypertarget{2-система-и-взаимодействие-частей}{}

Пользователь подтверждает свою личность, получает доступ к конфигурациям своей группы и создаёт набор задач. Конфигурация определяет состав набора; профильные домены формируют условия и проверяют решения. Домен сессий связывает задачи с пользователем, управляет сроком и собирает отдельные результаты в итог набора.

На системном уровне части рассматриваются через их ответственность и внешние обязательства. Для описания срока и завершения используется общая модель состояний. Самостоятельность частей не означает, что им назначены отдельные серверы или сервисы.

{\def\LTcaptype{none} % do not increment counter
\begin{longtable}[]{@{}
  >{\raggedright\arraybackslash}p{(\linewidth - 4\tabcolsep) * \real{0.2995}}
  >{\raggedright\arraybackslash}p{(\linewidth - 4\tabcolsep) * \real{0.4010}}
  >{\raggedright\arraybackslash}p{(\linewidth - 4\tabcolsep) * \real{0.2995}}@{}}
\toprule\noalign{}
\begin{minipage}[b]{\linewidth}\raggedright
Часть
\end{minipage} & \begin{minipage}[b]{\linewidth}\raggedright
Ответственность
\end{minipage} & \begin{minipage}[b]{\linewidth}\raggedright
Владелец
\end{minipage} \\
\midrule\noalign{}
\endhead
\bottomrule\noalign{}
\endlastfoot
Аутентификация и пользователи & Личность пользователя, роли, связь с группой & Кокорев \\
Сессии задач & Набор, ответы, срок и итог & Анохин \\
Ранцевые задачи & Условия и проверка ранцевых задач & Георгиев \\
Учебная бэкдор-атака & Предметный сценарий атаки & Минасов \\
Эль-Гамаль & Условия и проверка задач Эль-Гамаля & Портнов \\
Конфигурации и модерация & Состав доступных наборов и операции преподавателя & Самбьюи \\
\end{longtable}
}

\subsection{Пользовательский путь и границы ответственности}\label{ux43fux43eux43bux44cux437ux43eux432ux430ux442ux435ux43bux44cux441ux43aux438ux439-ux43fux443ux442ux44c-ux438-ux433ux440ux430ux43dux438ux446ux44b-ux43eux442ux432ux435ux442ux441ux442ux432ux435ux43dux43dux43eux441ux442ux438}

\hypertarget{пользовательский-путь-и-границы-ответственности}{}

Отправной точкой служит подтверждённая личность пользователя. Система использует её для выбора группы и проверки принадлежности сессии. Указание имени конфигурации или идентификатора задачи само по себе не даёт права выполнить операцию: эти сведения сопоставляются с данными текущего пользователя.

При создании набора система получает доступную конфигурацию и формирует условия предусмотренных ею задач. Результат создания связывает пользователя, сессию и её задачи. После выдачи набора проверка опирается на сохранённые условия; изменение конфигурации не должно изменять уже выданное задание. Поэтому выбор конфигурации и решение задачи относятся к разным этапам пользовательского пути.

При отправке ответа сначала проверяются условия допуска. Только допустимый по принадлежности, состоянию, сроку и формату ответ передаётся предметному проверяющему. Ошибка допуска сообщает о невозможности выполнить операцию, а предметный вердикт определяет результат решения. Это различие необходимо для применения правила единственной попытки.

Результат отдельной задачи и результат сессии имеют разный смысл. Неверный ответ завершает соответствующую задачу, но пока остаются другие незавершённые задачи, сессия сохраняет \texttt{NEW}. Итог сессии определяется после завершения всех её задач. Так пользовательские операции, предметные проверки и фоновое завершение используют общую модель результата.

Раздел сессий не задаёт математическую правильность решения и не устанавливает правила выдачи ролей. Эти обязательства принадлежат профильным владельцам. Для совместной работы здесь фиксируется необходимый внешний результат их операций: подтверждённый пользователь, доступный состав набора, сформированное условие и вердикт проверки.

\subsection{Переход от системного поведения к программным обязанностям}\label{ux43fux435ux440ux435ux445ux43eux434-ux43eux442-ux441ux438ux441ux442ux435ux43cux43dux43eux433ux43e-ux43fux43eux432ux435ux434ux435ux43dux438ux44f-ux43a-ux43fux440ux43eux433ux440ux430ux43cux43cux43dux44bux43c-ux43eux431ux44fux437ux430ux43dux43dux43eux441ux442ux44fux43c}

\hypertarget{переход-от-системного-поведения-к-программным-обязанностям}{}

Системное требование описывает наблюдаемое обязательство. Программные требования определяют, какие действия компонентов его обеспечивают. Например, SR-SES-011 сохраняет своевременность ответа при длительной серверной проверке. Из этого следуют две связанные обязанности: SWR-SES-007 фиксирует момент получения запроса, а SWR-SES-015 согласует его обработку с контролем просрочки.

Одна ссылка в таблице трассировки ещё не подтверждает полноту детализации. Для такого требования нужно проверить совместный результат: время поступления сохранено, фоновое завершение его учитывает, а вердикт после срока не подменяется просрочкой. Разделение по компонентам должно сохранять всё системное обязательство.

\section{Авторские разделы}\label{3-ux430ux432ux442ux43eux440ux441ux43aux438ux435-ux440ux430ux437ux434ux435ux43bux44b}

\hypertarget{3-авторские-разделы}{}

{\def\LTcaptype{none} % do not increment counter
\begin{longtable}[]{@{}
  >{\raggedright\arraybackslash}p{(\linewidth - 4\tabcolsep) * \real{0.3200}}
  >{\raggedright\arraybackslash}p{(\linewidth - 4\tabcolsep) * \real{0.3414}}
  >{\raggedright\arraybackslash}p{(\linewidth - 4\tabcolsep) * \real{0.3386}}@{}}
\toprule\noalign{}
\begin{minipage}[b]{\linewidth}\raggedright
Раздел
\end{minipage} & \begin{minipage}[b]{\linewidth}\raggedright
Где определён
\end{minipage} & \begin{minipage}[b]{\linewidth}\raggedright
Состояние
\end{minipage} \\
\midrule\noalign{}
\endhead
\bottomrule\noalign{}
\endlastfoot
Сессии задач & \hyperref[doc-05-ux442ux440ux435ux431ux43eux432ux430ux43dux438ux44f-ux43a-ux441ux435ux441ux441ux438ux44fux43c-ux437ux430ux434ux430ux447]{SR-SES-001--SR-SES-014} & Авторский раздел Анохина \\
Общесистемная интеграция & \hyperref[doc-05-ux442ux440ux435ux431ux43eux432ux430ux43dux438ux44f-ux43a-ux438ux43dux442ux435ux433ux440ux430ux446ux438ux438]{SR-SYS-001--SR-SYS-009} & Авторский раздел Анохина \\
Ранцевые задачи & \hyperref[doc-05-ux442ux440ux435ux431ux43eux432ux430ux43dux438ux44f-ux43a-ux440ux430ux43dux446ux435ux432ux44bux43c-ux437ux430ux434ux430ux447ux430ux43c]{SR-KNAP-001--SR-KNAP-010} & Авторский раздел Георгиева \\
Остальные домены, безопасность и эксплуатация & Авторские разделы владельцев & Самостоятельные части комплекта \\
\end{longtable}
}

Правила ID заданы в \href{../notes/methodology/\%D0\%9C\%D0\%B5\%D1\%82\%D0\%BE\%D0\%B4\%D0\%B8\%D0\%BA\%D0\%B0_\%D0\%B8\%D0\%B4\%D0\%B5\%D0\%BD\%D1\%82\%D0\%B8\%D1\%84\%D0\%B8\%D0\%BA\%D0\%B0\%D1\%82\%D0\%BE\%D1\%80\%D0\%BE\%D0\%B2.pdf\#nameddest=\%D0\%BC\%D0\%B5\%D1\%82\%D0\%BE\%D0\%B4\%D0\%B8\%D0\%BA\%D0\%B0-\%D0\%B8\%D0\%B4\%D0\%B5\%D0\%BD\%D1\%82\%D0\%B8\%D1\%84\%D0\%B8\%D0\%BA\%D0\%B0\%D1\%82\%D0\%BE\%D1\%80\%D0\%BE\%D0\%B2}{методике}; термины --- в \href{../notes/methodology/\%D0\%93\%D0\%BB\%D0\%BE\%D1\%81\%D1\%81\%D0\%B0\%D1\%80\%D0\%B8\%D0\%B9.pdf\#nameddest=\%D0\%B3\%D0\%BB\%D0\%BE\%D1\%81\%D1\%81\%D0\%B0\%D1\%80\%D0\%B8\%D0\%B9}{глоссарии}. Внешние обязанности определяются в авторских разделах один раз, а здесь связываются на уровне системы.

\section{Связь с ТЗ и программными требованиями}\label{4-ux441ux432ux44fux437ux44c-ux441-ux442ux437-ux438-ux43fux440ux43eux433ux440ux430ux43cux43cux43dux44bux43cux438-ux442ux440ux435ux431ux43eux432ux430ux43dux438ux44fux43cux438}

\hypertarget{4-связь-с-тз-и-программными-требованиями}{}

{\def\LTcaptype{none} % do not increment counter
\begin{longtable}[]{@{}
  >{\raggedright\arraybackslash}p{(\linewidth - 4\tabcolsep) * \real{0.2476}}
  >{\raggedright\arraybackslash}p{(\linewidth - 4\tabcolsep) * \real{0.3954}}
  >{\raggedright\arraybackslash}p{(\linewidth - 4\tabcolsep) * \real{0.3571}}@{}}
\toprule\noalign{}
\begin{minipage}[b]{\linewidth}\raggedright
Основание ТЗ
\end{minipage} & \begin{minipage}[b]{\linewidth}\raggedright
Системные обязательства
\end{minipage} & \begin{minipage}[b]{\linewidth}\raggedright
Программная детализация
\end{minipage} \\
\midrule\noalign{}
\endhead
\bottomrule\noalign{}
\endlastfoot
\texttt{[TS.\allowbreak{}REQ.\allowbreak{}FUNC.\allowbreak{}SESSION.\allowbreak{}CREATE]}, \texttt{[TS.\allowbreak{}REQ.\allowbreak{}FUNC.\allowbreak{}CONFIG.\allowbreak{}VIEW]} & Создание доступного набора и ограничение активных сессий & DOC-06.2: подготовка и создание \\
\texttt{[TS.\allowbreak{}REQ.\allowbreak{}FUNC.\allowbreak{}SESSION.\allowbreak{}SUBMIT]} & Принадлежность, выбор проверки, результат и итог сессии & DOC-06.2: обработка ответа \\
\texttt{[TS.\allowbreak{}REQ.\allowbreak{}FUNC.\allowbreak{}SESSION.\allowbreak{}TIME]} & Общий срок, просрочка и своевременность ответа & DOC-06.2: единый срок набора, учёт времени и планировщик \\
\texttt{{[}TS.REQ.INT{]}}, \texttt{[TS.\allowbreak{}REQ.\allowbreak{}REL.\allowbreak{}PERSIST]} & Контракты взаимодействия, общие данные и сохранность & DOC-06: общие требования; DOC-06.2: обязанности домена \\
\end{longtable}
}

Положение ТЗ о времени решения задач конкретизировано общим сроком набора: все задачи получают время создания сессии как начало отсчёта. Своевременность ответа определяется по получению полного запроса сервером; окончание проверки после срока допустимо в пределах установленного времени обработки.

Связи SR с определёнными SWR приведены в программных документах. Полные TC и общую матрицу выпускают владельцы DOC-10 и DOC-09; их готовность здесь не предполагается.

\section{Распределение ответственности}\label{5-ux440ux430ux441ux43fux440ux435ux434ux435ux43bux435ux43dux438ux435-ux43eux442ux432ux435ux442ux441ux442ux432ux435ux43dux43dux43eux441ux442ux438}

\hypertarget{5-распределение-ответственности}{}

{\def\LTcaptype{none} % do not increment counter
\begin{longtable}[]{@{}
  >{\raggedright\arraybackslash}p{(\linewidth - 4\tabcolsep) * \real{0.2597}}
  >{\raggedright\arraybackslash}p{(\linewidth - 4\tabcolsep) * \real{0.3285}}
  >{\raggedright\arraybackslash}p{(\linewidth - 4\tabcolsep) * \real{0.4118}}@{}}
\toprule\noalign{}
\begin{minipage}[b]{\linewidth}\raggedright
Владелец
\end{minipage} & \begin{minipage}[b]{\linewidth}\raggedright
Область документации
\end{minipage} & \begin{minipage}[b]{\linewidth}\raggedright
Граница взаимодействия с сессиями
\end{minipage} \\
\midrule\noalign{}
\endhead
\bottomrule\noalign{}
\endlastfoot
Кокорев & Аутентификация, пользователи и безопасность & Подтверждённый пользователь и роль доступа \\
Георгиев & Ранцевые криптосистемы & Генерация условия и предметная проверка ранцевого решения \\
Минасов & Бэкдор-сценарий и общий API & Внешние операции и единое представление ошибок \\
Портнов & Эль-Гамаль и трассировка & Генерация и проверка Эль-Гамаля, связи требований \\
Самбьюи & Конфигурации, хранение и инфраструктура & Доступный состав набора и сохранность данных \\
Анохин & Сессии и общесистемная интеграция & Срок, попытка, состояния и взаимодействие частей \\
\end{longtable}
}

Каждый автор определяет внутренние правила своей предметной области. Взаимодействие с сессиями выполняется по установленным контрактам. Изменение контракта отражается в системном требовании, программной детализации и описательной проверке; замена внутреннего алгоритма при сохранении контракта не требует переноса его описания в общий документ.

\section{История}\label{6-ux438ux441ux442ux43eux440ux438ux44f}

\hypertarget{6-история}{}

05.11.2026: подготовлена рабочая редакция второго спринта. Сохранены ID десяти ранее подготовленных Markdown-требований сессий; добавлены требование своевременности ответа, раздел интеграции и распределение ответственности. 07.11.2026: определены предел обработки, восстановление и форматы; расширены пояснения и подготовлены LaTeX и PDF.

Редакция 07.11.2026: развёрнуты пояснения каждого требования, добавлены примеры и точные ссылки на связанные разделы. Формулировки обязательств и идентификаторы сохранены.

\section{DOC-05. Требования к сессиям задач}\label{doc-05-ux442ux440ux435ux431ux43eux432ux430ux43dux438ux44f-ux43a-ux441ux435ux441ux441ux438ux44fux43c-ux437ux430ux434ux430ux447}

\hypertarget{doc-05-требования-к-сессиям-задач}{}

Раздел \hyperref[doc-05-ux441ux438ux441ux442ux435ux43cux43dux44bux435-ux442ux440ux435ux431ux43eux432ux430ux43dux438ux44f]{DOC-05}, версия 0.4 от 07.11.2026. Автор: Анохин Г.А. Все требования раздела имеют обязательный приоритет. Предметные основания: \href{../notes/domains/\%D0\%A1\%D0\%B5\%D1\%81\%D1\%81\%D0\%B8\%D0\%B8_\%D0\%B7\%D0\%B0\%D0\%B4\%D0\%B0\%D1\%87.pdf\#nameddest=\%D0\%BA\%D0\%B0\%D1\%80\%D1\%82\%D0\%BE\%D1\%87\%D0\%BA\%D0\%B0-\%D0\%B4\%D0\%BE\%D0\%BC\%D0\%B5\%D0\%BD\%D0\%B0-\%D1\%81\%D0\%B5\%D1\%81\%D1\%81\%D0\%B8\%D0\%B8-\%D0\%B7\%D0\%B0\%D0\%B4\%D0\%B0\%D1\%87}{карточка сессий} и указанные ниже требования \href{../02_\%D0\%A2\%D0\%B5\%D1\%85\%D0\%BD\%D0\%B8\%D1\%87\%D0\%B5\%D1\%81\%D0\%BA\%D0\%BE\%D0\%B5_\%D0\%B7\%D0\%B0\%D0\%B4\%D0\%B0\%D0\%BD\%D0\%B8\%D0\%B5/\%D0\%A2\%D0\%B5\%D1\%85\%D0\%BD\%D0\%B8\%D1\%87\%D0\%B5\%D1\%81\%D0\%BA\%D0\%BE\%D0\%B5_\%D0\%B7\%D0\%B0\%D0\%B4\%D0\%B0\%D0\%BD\%D0\%B8\%D0\%B5.pdf}{ТЗ}.

\subsection{SR-SES-001. Создание набора по конфигурации группы}\label{sr-ses-001-ux441ux43eux437ux434ux430ux43dux438ux435-ux43dux430ux431ux43eux440ux430-ux43fux43e-ux43aux43eux43dux444ux438ux433ux443ux440ux430ux446ux438ux438-ux433ux440ux443ux43fux43fux44b}

\hypertarget{sr-ses-001-создание-набора-по-конфигурации-группы}{}

Создание сессии превращает выбранную учебную конфигурацию в конкретный набор заданий одного пользователя. Конфигурация задаёт состав, а сессия хранит уже сформированные условия и результаты работы с ними. На этом переходе взаимодействуют проверка доступа, группа пользователя и генераторы криптографических задач; их обязанности разведены в \href{../07_\%D0\%9F\%D1\%80\%D0\%BE\%D0\%B3\%D1\%80\%D0\%B0\%D0\%BC\%D0\%BC\%D0\%BD\%D1\%8B\%D0\%B5_\%D1\%82\%D1\%80\%D0\%B5\%D0\%B1\%D0\%BE\%D0\%B2\%D0\%B0\%D0\%BD\%D0\%B8\%D1\%8F/DOC-06_\%D0\%9F\%D1\%80\%D0\%BE\%D0\%B3\%D1\%80\%D0\%B0\%D0\%BC\%D0\%BC\%D0\%BD\%D1\%8B\%D0\%B5_\%D1\%82\%D1\%80\%D0\%B5\%D0\%B1\%D0\%BE\%D0\%B2\%D0\%B0\%D0\%BD\%D0\%B8\%D1\%8F.pdf\#nameddest=2-\%D0\%BE\%D1\%82\%D0\%B2\%D0\%B5\%D1\%82\%D1\%81\%D1\%82\%D0\%B2\%D0\%B5\%D0\%BD\%D0\%BD\%D0\%BE\%D1\%81\%D1\%82\%D1\%8C-\%D0\%B8-\%D1\%81\%D0\%B2\%D1\%8F\%D0\%B7\%D0\%B8}{описании компонентов}. Пользователь выбирает доступный набор, а не передаёт серверу произвольный состав задач.

\textbf{Требование.} Система должна создавать сессию текущего пользователя по конфигурации, включённой для его группы, и возвращать данные созданной сессии с задачами, соответствующими составу конфигурации.

Если конфигурация включает ранцевую задачу и задачу Эль-Гамаля, в новой сессии сохраняются условия обоих заданий. Именно эти условия позднее используются при проверке, даже если настройки группы изменились. Создание считается законченным, когда полный набор сохранён и может быть возвращён пользователю. Программная последовательность раскрыта в \href{../07_\%D0\%9F\%D1\%80\%D0\%BE\%D0\%B3\%D1\%80\%D0\%B0\%D0\%BC\%D0\%BC\%D0\%BD\%D1\%8B\%D0\%B5_\%D1\%82\%D1\%80\%D0\%B5\%D0\%B1\%D0\%BE\%D0\%B2\%D0\%B0\%D0\%BD\%D0\%B8\%D1\%8F/DOC-06.2_\%D0\%A1\%D0\%B5\%D1\%81\%D1\%81\%D0\%B8\%D0\%B8_\%D0\%B7\%D0\%B0\%D0\%B4\%D0\%B0\%D1\%87_\%D0\%B8_\%D0\%BF\%D0\%BB\%D0\%B0\%D0\%BD\%D0\%B8\%D1\%80\%D0\%BE\%D0\%B2\%D1\%89\%D0\%B8\%D0\%BA.pdf\#nameddest=2-\%D1\%81\%D0\%BE\%D0\%B7\%D0\%B4\%D0\%B0\%D0\%BD\%D0\%B8\%D0\%B5-\%D1\%81\%D0\%B5\%D1\%81\%D1\%81\%D0\%B8\%D0\%B8}{разделе создания сессии}, а случай ошибки генерации --- в \href{../07_\%D0\%9F\%D1\%80\%D0\%BE\%D0\%B3\%D1\%80\%D0\%B0\%D0\%BC\%D0\%BC\%D0\%BD\%D1\%8B\%D0\%B5_\%D1\%82\%D1\%80\%D0\%B5\%D0\%B1\%D0\%BE\%D0\%B2\%D0\%B0\%D0\%BD\%D0\%B8\%D1\%8F/DOC-06.2_\%D0\%A1\%D0\%B5\%D1\%81\%D1\%81\%D0\%B8\%D0\%B8_\%D0\%B7\%D0\%B0\%D0\%B4\%D0\%B0\%D1\%87_\%D0\%B8_\%D0\%BF\%D0\%BB\%D0\%B0\%D0\%BD\%D0\%B8\%D1\%80\%D0\%BE\%D0\%B2\%D1\%89\%D0\%B8\%D0\%BA.pdf\#nameddest=swr-ses-005-\%D1\%81\%D0\%BE\%D1\%85\%D1\%80\%D0\%B0\%D0\%BD\%D0\%B5\%D0\%BD\%D0\%B8\%D0\%B5-\%D1\%81\%D0\%BE\%D0\%B7\%D0\%B4\%D0\%B0\%D0\%BD\%D0\%BD\%D0\%BE\%D0\%B3\%D0\%BE-\%D0\%BD\%D0\%B0\%D0\%B1\%D0\%BE\%D1\%80\%D0\%B0}{правиле сохранения полного набора}.

Основание: \texttt{[TS.\allowbreak{}REQ.\allowbreak{}FUNC.\allowbreak{}SESSION.\allowbreak{}CREATE]}, \texttt{[TS.\allowbreak{}REQ.\allowbreak{}FUNC.\allowbreak{}CONFIG.\allowbreak{}VIEW]}; сценарий создания набора.

Проверка: у пользователя есть группа с включённой конфигурацией и нет другой сессии \texttt{NEW}. Запрос создаёт его сессию в состоянии \texttt{NEW} с задачами выбранного состава. Недоступная конфигурация не должна приводить к созданию набора.

\subsection{SR-SES-002. Одна незавершённая сессия}\label{sr-ses-002-ux43eux434ux43dux430-ux43dux435ux437ux430ux432ux435ux440ux448ux451ux43dux43dux430ux44f-ux441ux435ux441ux441ux438ux44f}

\hypertarget{sr-ses-002-одна-незавершённая-сессия}{}

Ограничение относится к пользователю, а не к вкладке браузера или имени конфигурации. Без него можно было бы обходить общий срок, создавая несколько наборов и выбирая удобный для продолжения. Незавершённой считается сессия в состоянии \texttt{NEW}; наличие уже неверно решённой задачи не завершает весь набор, пока другие задачи ещё решаются. Это различие раскрывает \hyperref[sr-ses-007-ux438ux442ux43eux433-ux441ux435ux441ux441ux438ux438]{правило итога сессии}.

\textbf{Требование.} При наличии у пользователя сессии \texttt{NEW} система должна отклонять создание новой сессии.

Два одновременных нажатия на создание набора должны приводить к одному и тому же ограничению, что и два последовательных обращения. Предварительная проверка наличия сессии сама по себе не защищает от параллельных запросов: проверку и создание нужно рассматривать совместно. Если прежний набор уже просрочен, сначала применяется \href{../07_\%D0\%9F\%D1\%80\%D0\%BE\%D0\%B3\%D1\%80\%D0\%B0\%D0\%BC\%D0\%BC\%D0\%BD\%D1\%8B\%D0\%B5_\%D1\%82\%D1\%80\%D0\%B5\%D0\%B1\%D0\%BE\%D0\%B2\%D0\%B0\%D0\%BD\%D0\%B8\%D1\%8F/DOC-06.2_\%D0\%A1\%D0\%B5\%D1\%81\%D1\%81\%D0\%B8\%D0\%B8_\%D0\%B7\%D0\%B0\%D0\%B4\%D0\%B0\%D1\%87_\%D0\%B8_\%D0\%BF\%D0\%BB\%D0\%B0\%D0\%BD\%D0\%B8\%D1\%80\%D0\%BE\%D0\%B2\%D1\%89\%D0\%B8\%D0\%BA.pdf\#nameddest=4-\%D0\%BA\%D0\%BE\%D0\%BD\%D1\%82\%D1\%80\%D0\%BE\%D0\%BB\%D1\%8C-\%D1\%81\%D1\%80\%D0\%BE\%D0\%BA\%D0\%B0-\%D0\%B8-\%D0\%B8\%D1\%82\%D0\%BE\%D0\%B3}{порядок завершения по сроку}. Своевременный ответ в обработке сохраняет незавершённость набора до определения результата.

Основание: \texttt{[TS.\allowbreak{}REQ.\allowbreak{}FUNC.\allowbreak{}SESSION.\allowbreak{}CREATE]}; ограничение активных сессий в карточке домена.

Проверка: последовательные и одновременные запросы одного пользователя не создают две незавершённые сессии. Если прежняя сессия просрочена, её состояние сначала приводится в соответствие с правилами завершения.

\subsection{SR-SES-003. Принадлежность и состояние сессии}\label{sr-ses-003-ux43fux440ux438ux43dux430ux434ux43bux435ux436ux43dux43eux441ux442ux44c-ux438-ux441ux43eux441ux442ux43eux44fux43dux438ux435-ux441ux435ux441ux441ux438ux438}

\hypertarget{sr-ses-003-принадлежность-и-состояние-сессии}{}

Идентификатор задачи указывает объект операции, но не подтверждает право пользователя работать с ним. Проверка должна пройти по цепочке «задача --- сессия --- владелец» и сопоставить владельца с личностью из механизма аутентификации. Эти сведения нельзя брать из переданного клиентом имени или \texttt{userId}. Общий принцип установлен в \hyperref[sr-sys-001-ux435ux434ux438ux43dux430ux44f-ux43bux438ux447ux43dux43eux441ux442ux44c-ux43fux43eux43bux44cux437ux43eux432ux430ux442ux435ux43bux44f]{требовании единой личности пользователя}.

\textbf{Требование.} Система должна допускать ответ к предметной проверке только для задачи собственной сессии текущего пользователя, если сессия находится в состоянии \texttt{NEW}.

Для собственного набора дополнительно проверяется состояние сессии. Завершённая сессия сохраняет результаты, но уже не принимает решения. Обращение к чужому заданию не должно запускать математический проверяющий и изменять данные. Внешний отказ задан в \href{../07_\%D0\%9F\%D1\%80\%D0\%BE\%D0\%B3\%D1\%80\%D0\%B0\%D0\%BC\%D0\%BC\%D0\%BD\%D1\%8B\%D0\%B5_\%D1\%82\%D1\%80\%D0\%B5\%D0\%B1\%D0\%BE\%D0\%B2\%D0\%B0\%D0\%BD\%D0\%B8\%D1\%8F/DOC-06.2_\%D0\%A1\%D0\%B5\%D1\%81\%D1\%81\%D0\%B8\%D0\%B8_\%D0\%B7\%D0\%B0\%D0\%B4\%D0\%B0\%D1\%87_\%D0\%B8_\%D0\%BF\%D0\%BB\%D0\%B0\%D0\%BD\%D0\%B8\%D1\%80\%D0\%BE\%D0\%B2\%D1\%89\%D0\%B8\%D0\%BA.pdf\#nameddest=\%D0\%BA\%D0\%BE\%D0\%B4\%D1\%8B-\%D0\%B8-\%D0\%BF\%D1\%80\%D0\%B5\%D0\%B4\%D1\%81\%D1\%82\%D0\%B0\%D0\%B2\%D0\%BB\%D0\%B5\%D0\%BD\%D0\%B8\%D0\%B5-\%D0\%BE\%D1\%88\%D0\%B8\%D0\%B1\%D0\%BE\%D0\%BA}{контракте ошибок}; доступ к результатам собственной завершённой сессии сохраняется через \hyperref[sr-ses-014-ux43fux43eux43bux443ux447ux435ux43dux438ux435-ux441ux43eux445ux440ux430ux43dux451ux43dux43dux43eux433ux43e-ux440ux435ux437ux443ux43bux44cux442ux430ux442ux430]{операцию чтения}.

Основание: \texttt{[TS.\allowbreak{}REQ.\allowbreak{}FUNC.\allowbreak{}SESSION.\allowbreak{}SUBMIT]}; сценарий отправки ответа.

Проверка: обращение к задаче чужой или завершённой сессии отклоняется до предметной проверки; её данные и результаты не меняются. Срок и состояние самой задачи дополнительно проверяются по SR-SES-006, SR-SES-008 и SR-SES-011.

\subsection{SR-SES-004. Выбор предметной проверки}\label{sr-ses-004-ux432ux44bux431ux43eux440-ux43fux440ux435ux434ux43cux435ux442ux43dux43eux439-ux43fux440ux43eux432ux435ux440ux43aux438}

\hypertarget{sr-ses-004-выбор-предметной-проверки}{}

Домен сессий управляет жизненным циклом задачи, но не решает, верно ли выполнен криптографический расчёт. Для этого существуют предметные проверяющие. Сессия передаёт им сохранённое условие и ответ пользователя, а получает вердикт либо категорию отказа. Такой внешний контракт позволяет описывать криптографические домены как самостоятельные части без переноса их алгоритмов в системные требования.

\textbf{Требование.} Для допустимого ответа система должна выбирать проверку по типу сохранённой задачи и передавать ей условие этой задачи и полученный ответ.

Ответ на задачу Эль-Гамаля не должен попасть в проверяющий ранцевой задачи даже тогда, когда клиент указал в теле иной тип. Источник выбора --- сохранённая задача. Содержимое и типы полей раскрыты в \href{../07_\%D0\%9F\%D1\%80\%D0\%BE\%D0\%B3\%D1\%80\%D0\%B0\%D0\%BC\%D0\%BC\%D0\%BD\%D1\%8B\%D0\%B5_\%D1\%82\%D1\%80\%D0\%B5\%D0\%B1\%D0\%BE\%D0\%B2\%D0\%B0\%D0\%BD\%D0\%B8\%D1\%8F/DOC-06.2_\%D0\%A1\%D0\%B5\%D1\%81\%D1\%81\%D0\%B8\%D0\%B8_\%D0\%B7\%D0\%B0\%D0\%B4\%D0\%B0\%D1\%87_\%D0\%B8_\%D0\%BF\%D0\%BB\%D0\%B0\%D0\%BD\%D0\%B8\%D1\%80\%D0\%BE\%D0\%B2\%D1\%89\%D0\%B8\%D0\%BA.pdf\#nameddest=\%D0\%BF\%D1\%80\%D0\%B5\%D0\%B4\%D0\%BC\%D0\%B5\%D1\%82\%D0\%BD\%D1\%8B\%D0\%B9-\%D0\%BE\%D1\%82\%D0\%B2\%D0\%B5\%D1\%82}{предметном контракте ответа}, а порядок вызова --- в \href{../07_\%D0\%9F\%D1\%80\%D0\%BE\%D0\%B3\%D1\%80\%D0\%B0\%D0\%BC\%D0\%BC\%D0\%BD\%D1\%8B\%D0\%B5_\%D1\%82\%D1\%80\%D0\%B5\%D0\%B1\%D0\%BE\%D0\%B2\%D0\%B0\%D0\%BD\%D0\%B8\%D1\%8F/DOC-06.2_\%D0\%A1\%D0\%B5\%D1\%81\%D1\%81\%D0\%B8\%D0\%B8_\%D0\%B7\%D0\%B0\%D0\%B4\%D0\%B0\%D1\%87_\%D0\%B8_\%D0\%BF\%D0\%BB\%D0\%B0\%D0\%BD\%D0\%B8\%D1\%80\%D0\%BE\%D0\%B2\%D1\%89\%D0\%B8\%D0\%BA.pdf\#nameddest=swr-ses-010-\%D0\%BF\%D0\%B5\%D1\%80\%D0\%B5\%D0\%B4\%D0\%B0\%D1\%87\%D0\%B0-\%D0\%BF\%D1\%80\%D0\%BE\%D1\%84\%D0\%B8\%D0\%BB\%D1\%8C\%D0\%BD\%D0\%BE\%D0\%B9-\%D0\%BF\%D1\%80\%D0\%BE\%D0\%B2\%D0\%B5\%D1\%80\%D0\%BA\%D0\%B5}{программном требовании передачи проверке}. Недоступность нужного обработчика означает технический отказ, а не неверное решение.

Основание: \texttt{[TS.\allowbreak{}REQ.\allowbreak{}FUNC.\allowbreak{}SESSION.\allowbreak{}SUBMIT]}; ответственность криптографических доменов в карточке сессий.

Проверка: для каждого поддерживаемого типа вызывается соответствующая предметная проверка. Тип, указанный пользователем в ответе, не должен подменять тип выбранной задачи.

\subsection{SR-SES-005. Результат задачи}\label{sr-ses-005-ux440ux435ux437ux443ux43bux44cux442ux430ux442-ux437ux430ux434ux430ux447ux438}

\hypertarget{sr-ses-005-результат-задачи}{}

Получение вердикта от математического проверяющего ещё не завершает пользовательскую операцию. Результат должен стать частью данных задачи и согласоваться с состоянием её сессии. Иначе пользователь мог бы увидеть успешный ответ, а после перезагрузки --- другое состояние. Граница подтверждения определяется \hyperref[sr-sys-005-ux441ux43eux433ux43bux430ux441ux43eux432ux430ux43dux43dux43eux441ux442ux44c-ux440ux435ux437ux443ux43bux44cux442ux430ux442ux430-ux438-ux441ux43eux441ux442ux43eux44fux43dux438ux44f]{требованием согласованности результата и состояния}.

\textbf{Требование.} После завершённой предметной проверки система должна сохранить задаче \texttt{SOLVED} при верном ответе либо \texttt{FAILED} при неверном и вернуть результат, относящийся к этой задаче.

При верном ответе сохраняются \texttt{SOLVED} и положительный предметный результат. При неверном --- \texttt{FAILED} и отрицательный результат с пояснением. Дальше пересчитывается итог сессии: отдельная ошибка не закрывает оставшиеся задачи. Повторное чтение возвращает уже сохранённый результат, не обращаясь к проверяющему. Различия полей результата и технического отказа приведены в \href{../07_\%D0\%9F\%D1\%80\%D0\%BE\%D0\%B3\%D1\%80\%D0\%B0\%D0\%BC\%D0\%BC\%D0\%BD\%D1\%8B\%D0\%B5_\%D1\%82\%D1\%80\%D0\%B5\%D0\%B1\%D0\%BE\%D0\%B2\%D0\%B0\%D0\%BD\%D0\%B8\%D1\%8F/DOC-06.2_\%D0\%A1\%D0\%B5\%D1\%81\%D1\%81\%D0\%B8\%D0\%B8_\%D0\%B7\%D0\%B0\%D0\%B4\%D0\%B0\%D1\%87_\%D0\%B8_\%D0\%BF\%D0\%BB\%D0\%B0\%D0\%BD\%D0\%B8\%D1\%80\%D0\%BE\%D0\%B2\%D1\%89\%D0\%B8\%D0\%BA.pdf\#nameddest=\%D1\%80\%D0\%B5\%D0\%B7\%D1\%83\%D0\%BB\%D1\%8C\%D1\%82\%D0\%B0\%D1\%82-\%D0\%B8-\%D0\%BE\%D1\%82\%D0\%BA\%D0\%B0\%D0\%B7\%D1\%8B}{контракте результата}.

Основание: \texttt{[TS.\allowbreak{}REQ.\allowbreak{}FUNC.\allowbreak{}SESSION.\allowbreak{}SUBMIT]}; сценарий проверки ответа.

Проверка: результат содержит идентификатор и тип проверенной задачи; сохранённое состояние соответствует вердикту. Технический сбой без вердикта обрабатывается по SR-SES-010.

\subsection{SR-SES-006. Общий срок сессии}\label{sr-ses-006-ux43eux431ux449ux438ux439-ux441ux440ux43eux43a-ux441ux435ux441ux441ux438ux438}

\hypertarget{sr-ses-006-общий-срок-сессии}{}

Срок относится ко всему набору. Он отсчитывается от одного серверного времени создания сессии, поэтому получение новых данных, перезагрузка страницы или повтор запроса не создают новый интервал решения. Клиент может отображать оставшееся время, но допуск определяет сервер. Обозначения \texttt{t0}, \texttt{D} и \texttt{t\_in} введены в \href{../notes/domains/\%D0\%A1\%D0\%B5\%D1\%81\%D1\%81\%D0\%B8\%D0\%B8_\%D0\%B7\%D0\%B0\%D0\%B4\%D0\%B0\%D1\%87.pdf\#nameddest=\%D1\%81\%D1\%80\%D0\%BE\%D0\%BA-\%D0\%B8-\%D0\%BF\%D0\%BE\%D0\%BF\%D1\%8B\%D1\%82\%D0\%BA\%D0\%B0}{карточке домена}.

\textbf{Требование.} Система должна ограничивать решение всех задач единым сроком: 20 минут от серверного времени создания сессии. По достижении срока незавершённые задачи без своевременного ответа в обработке должны завершаться с \texttt{FAILED}, а ранее определённые результаты --- сохраняться.

Нужно различать конец времени решения и момент фоновой записи просрочки. Планировщик может обновить хранилище позднее, но это не разрешает новые поздние ответы. При этом уже полученный своевременный ответ защищён от назначения просрочки, пока обрабатывается в допустимых пределах. Его поведение определяет \hyperref[sr-ses-011-ux43cux43eux43cux435ux43dux442-ux43fux440ux438ux451ux43cux430-ux43eux442ux432ux435ux442ux430]{правило момента приёма}, а окончание ожидания --- \hyperref[sr-ses-012-ux43eux433ux440ux430ux43dux438ux447ux435ux43dux43dux430ux44f-ux434ux43bux438ux442ux435ux43bux44cux43dux43eux441ux442ux44c-ux43eux431ux440ux430ux431ux43eux442ux43aux438]{предел обработки}.

Основание: сценарий срока в карточке домена; \texttt{[TS.\allowbreak{}REQ.\allowbreak{}FUNC.\allowbreak{}SESSION.\allowbreak{}TIME]}. Все задачи набора получают время создания сессии как начало отсчёта своего срока.

Проверка: задачи одной сессии имеют один срок. Обновление страницы и повторы запроса его не меняют. После срока необработанные задачи получают \texttt{FAILED}; своевременно полученный ответ обрабатывается по SR-SES-011. Задержка фоновой обработки не разрешает новые ответы после срока.

\subsection{SR-SES-007. Итог сессии}\label{sr-ses-007-ux438ux442ux43eux433-ux441ux435ux441ux441ux438ux438}

\hypertarget{sr-ses-007-итог-сессии}{}

Состояние сессии подводит итог всему набору и не копируется из последней обработанной задачи. Для этого система рассматривает состояния всех задач. Пока хотя бы одна задача не имеет итогового состояния, набор остаётся незавершённым. Остальные задачи сохраняют свои состояния и, если предметная проверка состоялась, зафиксированные вердикты. У завершённой по просрочке задачи состояние \texttt{FAILED}, но предметный вердикт отсутствует.

\textbf{Требование.} После завершения всех задач система должна устанавливать сессии \texttt{SOLVED}, если все задачи имеют \texttt{SOLVED}, и \texttt{FAILED}, если хотя бы одна задача имеет \texttt{FAILED}.

Если ранцевая задача решена неверно, а ответ на Эль-Гамаля ещё проверяется, итог сессии пока \texttt{NEW}. После завершения этой проверки набор получит \texttt{FAILED} из-за ранцевой задачи, даже при верном втором ответе. Успешный итог возможен только при успехе каждой задачи. Одинаковое агрегирование после ответа и планировщика закреплено в \href{../07_\%D0\%9F\%D1\%80\%D0\%BE\%D0\%B3\%D1\%80\%D0\%B0\%D0\%BC\%D0\%BC\%D0\%BD\%D1\%8B\%D0\%B5_\%D1\%82\%D1\%80\%D0\%B5\%D0\%B1\%D0\%BE\%D0\%B2\%D0\%B0\%D0\%BD\%D0\%B8\%D1\%8F/DOC-06.2_\%D0\%A1\%D0\%B5\%D1\%81\%D1\%81\%D0\%B8\%D0\%B8_\%D0\%B7\%D0\%B0\%D0\%B4\%D0\%B0\%D1\%87_\%D0\%B8_\%D0\%BF\%D0\%BB\%D0\%B0\%D0\%BD\%D0\%B8\%D1\%80\%D0\%BE\%D0\%B2\%D1\%89\%D0\%B8\%D0\%BA.pdf\#nameddest=swr-ses-018-\%D0\%BE\%D0\%B1\%D1\%89\%D0\%B0\%D1\%8F-\%D1\%84\%D1\%83\%D0\%BD\%D0\%BA\%D1\%86\%D0\%B8\%D1\%8F-\%D0\%B8\%D1\%82\%D0\%BE\%D0\%B3\%D0\%B0-\%D1\%81\%D0\%B5\%D1\%81\%D1\%81\%D0\%B8\%D0\%B8}{общей функции итога}.

Основание: \texttt{[TS.\allowbreak{}REQ.\allowbreak{}FUNC.\allowbreak{}SESSION.\allowbreak{}SUBMIT]}; правило завершения в карточке домена.

Проверка: рассматриваются все верные ответы, наличие неверного ответа и наличие просрочки. Пока есть задача без итогового состояния, в том числе своевременный ответ в обработке, сессия остаётся \texttt{NEW}.

\subsection{SR-SES-008. Одна предметная попытка}\label{sr-ses-008-ux43eux434ux43dux430-ux43fux440ux435ux434ux43cux435ux442ux43dux430ux44f-ux43fux43eux43fux44bux442ux43aux430}

\hypertarget{sr-ses-008-одна-предметная-попытка}{}

Попытка означает предметный результат, зафиксированный для задачи. Она не равна числу HTTP-запросов: запрос может быть отклонён до проверки, а вычисление --- прерваться без сохранения вердикта. Связь попытки с фиксацией нужна и для восстановления после сбоя. Термины используются одинаково в карточке, системных требованиях и \href{../07_\%D0\%9F\%D1\%80\%D0\%BE\%D0\%B3\%D1\%80\%D0\%B0\%D0\%BC\%D0\%BC\%D0\%BD\%D1\%8B\%D0\%B5_\%D1\%82\%D1\%80\%D0\%B5\%D0\%B1\%D0\%BE\%D0\%B2\%D0\%B0\%D0\%BD\%D0\%B8\%D1\%8F/DOC-06.2_\%D0\%A1\%D0\%B5\%D1\%81\%D1\%81\%D0\%B8\%D0\%B8_\%D0\%B7\%D0\%B0\%D0\%B4\%D0\%B0\%D1\%87_\%D0\%B8_\%D0\%BF\%D0\%BB\%D0\%B0\%D0\%BD\%D0\%B8\%D1\%80\%D0\%BE\%D0\%B2\%D1\%89\%D0\%B8\%D0\%BA.pdf\#nameddest=\%D0\%BF\%D1\%80\%D0\%B8\%D0\%BD\%D1\%8F\%D1\%82\%D0\%B0\%D1\%8F-\%D0\%BE\%D0\%B1\%D1\%80\%D0\%B0\%D0\%B1\%D0\%BE\%D1\%82\%D0\%BA\%D0\%B0-\%D0\%B8-\%D0\%B2\%D0\%BE\%D1\%81\%D1\%81\%D1\%82\%D0\%B0\%D0\%BD\%D0\%BE\%D0\%B2\%D0\%BB\%D0\%B5\%D0\%BD\%D0\%B8\%D0\%B5}{контракте принятой операции}.

\textbf{Требование.} Система должна фиксировать не более одного предметного вердикта для каждой задачи; после получения итогового состояния повторный ответ не должен его изменять.

Повтор после сохранённого вердикта не даёт ещё одной возможности изменить результат. Одновременная отправка тоже не должна создать два независимых итога. Если же первая операция завершилась ошибкой формата или техническим отказом без вердикта, действует правило нового обращения в пределах общего срока. Эти случаи описаны в \hyperref[sr-ses-009-ux43eux448ux438ux431ux43aux430-ux444ux43eux440ux43cux430ux442ux430]{обработке формата} и \hyperref[sr-ses-010-ux442ux435ux445ux43dux438ux447ux435ux441ux43aux438ux439-ux441ux431ux43eux439-ux43fux440ux43eux432ux435ux440ux43aux438]{обработке технического сбоя}.

Основание: правило попытки в карточке домена.

Проверка: повторные и одновременные отправки не приводят к двум предметным вердиктам. Ответ после \texttt{SOLVED} или \texttt{FAILED} отклоняется, результат сохраняется. Исключения без расходования попытки определены в SR-SES-009 и SR-SES-010.

\subsection{SR-SES-009. Ошибка формата}\label{sr-ses-009-ux43eux448ux438ux431ux43aux430-ux444ux43eux440ux43cux430ux442ux430}

\hypertarget{sr-ses-009-ошибка-формата}{}

Проверка формата отвечает на вопрос, можно ли передать данные предметному алгоритму. Она проверяет наличие тела, обязательных полей и ожидаемые типы данных. Правильно оформленный, но математически неверный расчёт относится уже к предметной проверке. Структура обоих поддерживаемых видов решения определена в \href{../07_\%D0\%9F\%D1\%80\%D0\%BE\%D0\%B3\%D1\%80\%D0\%B0\%D0\%BC\%D0\%BC\%D0\%BD\%D1\%8B\%D0\%B5_\%D1\%82\%D1\%80\%D0\%B5\%D0\%B1\%D0\%BE\%D0\%B2\%D0\%B0\%D0\%BD\%D0\%B8\%D1\%8F/DOC-06.2_\%D0\%A1\%D0\%B5\%D1\%81\%D1\%81\%D0\%B8\%D0\%B8_\%D0\%B7\%D0\%B0\%D0\%B4\%D0\%B0\%D1\%87_\%D0\%B8_\%D0\%BF\%D0\%BB\%D0\%B0\%D0\%BD\%D0\%B8\%D1\%80\%D0\%BE\%D0\%B2\%D1\%89\%D0\%B8\%D0\%BA.pdf\#nameddest=\%D0\%BF\%D1\%80\%D0\%B5\%D0\%B4\%D0\%BC\%D0\%B5\%D1\%82\%D0\%BD\%D1\%8B\%D0\%B9-\%D0\%BE\%D1\%82\%D0\%B2\%D0\%B5\%D1\%82}{контракте ответа}.

\textbf{Требование.} Система должна отклонять запрос с некорректным форматом или отсутствующим обязательным полем до предметной проверки, сообщать причину и сохранять неиспользованную попытку.

Например, отсутствующее обязательное поле не должно автоматически давать задаче \texttt{FAILED} как неверно решённой. Клиент получает категорию ошибки и может исправить запрос до общего срока. Время не начинается заново: исправленный запрос имеет собственный момент поступления. Если он получен после \texttt{D}, применяется \hyperref[sr-ses-011-ux43cux43eux43cux435ux43dux442-ux43fux440ux438ux451ux43cux430-ux43eux442ux432ux435ux442ux430]{правило своевременности}, независимо от времени первого обращения.

Основание: правило обработки ошибок в карточке домена.

Проверка: неверно оформленный ответ не получает предметного вердикта. Исправленный ответ допускается, если новый запрос получен до срока и выполнены остальные условия. Состояние \texttt{NEW} сохраняется, пока задача не завершена по другому основанию; сама ошибка срока не продлевает.

\subsection{SR-SES-010. Технический сбой проверки}\label{sr-ses-010-ux442ux435ux445ux43dux438ux447ux435ux441ux43aux438ux439-ux441ux431ux43eux439-ux43fux440ux43eux432ux435ux440ux43aux438}

\hypertarget{sr-ses-010-технический-сбой-проверки}{}

Технический отказ означает, что система не получила и не зафиксировала предметный результат. Причиной может стать недоступность проверяющего, прерванное вычисление или невозможность сохранить результат. Такой отказ не устанавливает правильность ответа пользователя и не должен маскироваться под отрицательный вердикт. Внешние категории разделены в \href{../07_\%D0\%9F\%D1\%80\%D0\%BE\%D0\%B3\%D1\%80\%D0\%B0\%D0\%BC\%D0\%BC\%D0\%BD\%D1\%8B\%D0\%B5_\%D1\%82\%D1\%80\%D0\%B5\%D0\%B1\%D0\%BE\%D0\%B2\%D0\%B0\%D0\%BD\%D0\%B8\%D1\%8F/DOC-06.2_\%D0\%A1\%D0\%B5\%D1\%81\%D1\%81\%D0\%B8\%D0\%B8_\%D0\%B7\%D0\%B0\%D0\%B4\%D0\%B0\%D1\%87_\%D0\%B8_\%D0\%BF\%D0\%BB\%D0\%B0\%D0\%BD\%D0\%B8\%D1\%80\%D0\%BE\%D0\%B2\%D1\%89\%D0\%B8\%D0\%BA.pdf\#nameddest=\%D0\%BA\%D0\%BE\%D0\%B4\%D1\%8B-\%D0\%B8-\%D0\%BF\%D1\%80\%D0\%B5\%D0\%B4\%D1\%81\%D1\%82\%D0\%B0\%D0\%B2\%D0\%BB\%D0\%B5\%D0\%BD\%D0\%B8\%D0\%B5-\%D0\%BE\%D1\%88\%D0\%B8\%D0\%B1\%D0\%BE\%D0\%BA}{таблице ошибок}.

\textbf{Требование.} Если проверка не завершилась вердиктом из-за технического сбоя, система должна сообщить о сбое и сохранить неиспользованную предметную попытку.

Сохранённая попытка не даёт бессрочного права продолжать набор. До \texttt{D} пользователь может отправить новый допустимый ответ; после \texttt{D} такой запрос будет поздним. Если прежний своевременный запрос уже исчерпал \hyperref[sr-ses-012-ux43eux433ux440ux430ux43dux438ux447ux435ux43dux43dux430ux44f-ux434ux43bux438ux442ux435ux43bux44cux43dux43eux441ux442ux44c-ux43eux431ux440ux430ux431ux43eux442ux43aux438]{предел обработки}, его поздний результат не записывается. Порядок возобновления после перезапуска раскрыт в \hyperref[sr-ses-013-ux432ux43eux441ux441ux442ux430ux43dux43eux432ux43bux435ux43dux438ux435-ux43fux440ux438ux43dux44fux442ux43eux433ux43e-ux43eux442ux432ux435ux442ux430]{требовании восстановления}.

Основание: правило технического сбоя в карточке домена.

Проверка: до срока возможен новый допустимый запрос. Если сбой своевременной проверки обнаружен после срока, новый ответ уже не принимается, а незавершённая задача закрывается по просрочке. Ошибка сервера не считается неверным предметным ответом.

\subsection{SR-SES-011. Момент приёма ответа}\label{sr-ses-011-ux43cux43eux43cux435ux43dux442-ux43fux440ux438ux451ux43cux430-ux43eux442ux432ux435ux442ux430}

\hypertarget{sr-ses-011-момент-приёма-ответа}{}

Пользователь не управляет длительностью серверной очереди и математической проверки. Поэтому своевременность связывается с получением полного запроса сервером, а не с окончанием вычисления. Отметка времени на устройстве пользователя для такого решения не подходит: система использует собственное \texttt{t\_in} и сравнивает его с общим сроком.

\textbf{Требование.} Система должна считать ответ своевременным, если полный запрос получен сервером строго до общего срока сессии, и сохранять право на его проверку после наступления этого срока.

Ответ, полученный до \texttt{D}, может завершиться вердиктом после \texttt{D}, если результат зафиксирован до \texttt{t\_in\ +\ 60\ секунд}. После этого предела поздний вердикт не сохраняется; действует \hyperref[sr-ses-012-ux43eux433ux440ux430ux43dux438ux447ux435ux43dux43dux430ux44f-ux434ux43bux438ux442ux435ux43bux44cux43dux43eux441ux442ux44c-ux43eux431ux440ux430ux431ux43eux442ux43aux438]{ограничение обработки SR-SES-012}. Запрос, полученный точно в \texttt{D}, уже не допускается. Эти события нельзя объединять в одну проверку состояния \texttt{NEW}: запись фоновой просрочки может ещё не выполниться. Регистрацию времени и защиту от планировщика связывают \href{../07_\%D0\%9F\%D1\%80\%D0\%BE\%D0\%B3\%D1\%80\%D0\%B0\%D0\%BC\%D0\%BC\%D0\%BD\%D1\%8B\%D0\%B5_\%D1\%82\%D1\%80\%D0\%B5\%D0\%B1\%D0\%BE\%D0\%B2\%D0\%B0\%D0\%BD\%D0\%B8\%D1\%8F/DOC-06.2_\%D0\%A1\%D0\%B5\%D1\%81\%D1\%81\%D0\%B8\%D0\%B8_\%D0\%B7\%D0\%B0\%D0\%B4\%D0\%B0\%D1\%87_\%D0\%B8_\%D0\%BF\%D0\%BB\%D0\%B0\%D0\%BD\%D0\%B8\%D1\%80\%D0\%BE\%D0\%B2\%D1\%89\%D0\%B8\%D0\%BA.pdf\#nameddest=swr-ses-007-\%D1\%80\%D0\%B5\%D0\%B3\%D0\%B8\%D1\%81\%D1\%82\%D1\%80\%D0\%B0\%D1\%86\%D0\%B8\%D1\%8F-\%D0\%B2\%D1\%80\%D0\%B5\%D0\%BC\%D0\%B5\%D0\%BD\%D0\%B8-\%D0\%BF\%D0\%BE\%D1\%81\%D1\%82\%D1\%83\%D0\%BF\%D0\%BB\%D0\%B5\%D0\%BD\%D0\%B8\%D1\%8F}{учёт поступления} и \href{../07_\%D0\%9F\%D1\%80\%D0\%BE\%D0\%B3\%D1\%80\%D0\%B0\%D0\%BC\%D0\%BC\%D0\%BD\%D1\%8B\%D0\%B5_\%D1\%82\%D1\%80\%D0\%B5\%D0\%B1\%D0\%BE\%D0\%B2\%D0\%B0\%D0\%BD\%D0\%B8\%D1\%8F/DOC-06.2_\%D0\%A1\%D0\%B5\%D1\%81\%D1\%81\%D0\%B8\%D0\%B8_\%D0\%B7\%D0\%B0\%D0\%B4\%D0\%B0\%D1\%87_\%D0\%B8_\%D0\%BF\%D0\%BB\%D0\%B0\%D0\%BD\%D0\%B8\%D1\%80\%D0\%BE\%D0\%B2\%D1\%89\%D0\%B8\%D0\%BA.pdf\#nameddest=swr-ses-015-\%D0\%B7\%D0\%B0\%D0\%B2\%D0\%B5\%D1\%80\%D1\%88\%D0\%B5\%D0\%BD\%D0\%B8\%D0\%B5-\%D1\%81\%D0\%B2\%D0\%BE\%D0\%B5\%D0\%B2\%D1\%80\%D0\%B5\%D0\%BC\%D0\%B5\%D0\%BD\%D0\%BD\%D0\%BE\%D0\%B9-\%D0\%BF\%D1\%80\%D0\%BE\%D0\%B2\%D0\%B5\%D1\%80\%D0\%BA\%D0\%B8}{завершение своевременной проверки}.

Основание: раздел «Срок и попытка» карточки домена.

Проверка: ответ получен за секунду до срока, а вердикт готов через секунду после него --- вердикт принимается. Запрос, полученный точно в срок или позднее, отклоняется. Зафиксированное сервером время получения, а не время окончания проверки или часы пользователя, определяет своевременность.

Своевременность не отменяет проверок принадлежности, состояния и формата. Ожидание работы сервера не создаёт дополнительной попытки. Программные связи всех требований приведены в \href{../07_\%D0\%9F\%D1\%80\%D0\%BE\%D0\%B3\%D1\%80\%D0\%B0\%D0\%BC\%D0\%BC\%D0\%BD\%D1\%8B\%D0\%B5_\%D1\%82\%D1\%80\%D0\%B5\%D0\%B1\%D0\%BE\%D0\%B2\%D0\%B0\%D0\%BD\%D0\%B8\%D1\%8F/DOC-06.2_\%D0\%A1\%D0\%B5\%D1\%81\%D1\%81\%D0\%B8\%D0\%B8_\%D0\%B7\%D0\%B0\%D0\%B4\%D0\%B0\%D1\%87_\%D0\%B8_\%D0\%BF\%D0\%BB\%D0\%B0\%D0\%BD\%D0\%B8\%D1\%80\%D0\%BE\%D0\%B2\%D1\%89\%D0\%B8\%D0\%BA.pdf\#nameddest=6-\%D1\%82\%D1\%80\%D0\%B0\%D1\%81\%D1\%81\%D0\%B8\%D1\%80\%D0\%BE\%D0\%B2\%D0\%BA\%D0\%B0-\%D0\%B8-\%D0\%BE\%D0\%BF\%D0\%B8\%D1\%81\%D0\%B0\%D1\%82\%D0\%B5\%D0\%BB\%D1\%8C\%D0\%BD\%D1\%8B\%D0\%B5-\%D0\%BF\%D1\%80\%D0\%BE\%D0\%B2\%D0\%B5\%D1\%80\%D0\%BA\%D0\%B8}{DOC-06.2}.

\subsection{SR-SES-012. Ограниченная длительность обработки}\label{sr-ses-012-ux43eux433ux440ux430ux43dux438ux447ux435ux43dux43dux430ux44f-ux434ux43bux438ux442ux435ux43bux44cux43dux43eux441ux442ux44c-ux43eux431ux440ux430ux431ux43eux442ux43aux438}

\hypertarget{sr-ses-012-ограниченная-длительность-обработки}{}

Право завершить своевременную проверку после общего срока не означает неограниченное ожидание. У принятой операции есть собственный предел, который включает очередь, вычисление и восстановление. Система должна иметь определённый исход и тогда, когда проверяющий завис или не смог вернуть ответ. Этот предел отличается от периода запуска фонового планировщика.

\textbf{Требование.} Система должна фиксировать результат обработки принятого ответа до момента «получение запроса + 60 секунд». Начиная с этого момента операция без сохранённого вердикта завершается техническим отказом без расходования предметной попытки.

Если результат не зафиксирован до \texttt{t\_in\ +\ 60\ секунд}, операция прекращается с техническим отказом. Вычисление, закончившееся позже, уже не меняет задачу. Далее учитывается общий срок: до \texttt{D} возможен новый запрос, после \texttt{D} задача закрывается по просрочке. Обе временные границы сопоставлены в \href{../07_\%D0\%9F\%D1\%80\%D0\%BE\%D0\%B3\%D1\%80\%D0\%B0\%D0\%BC\%D0\%BC\%D0\%BD\%D1\%8B\%D0\%B5_\%D1\%82\%D1\%80\%D0\%B5\%D0\%B1\%D0\%BE\%D0\%B2\%D0\%B0\%D0\%BD\%D0\%B8\%D1\%8F/DOC-06.2_\%D0\%A1\%D0\%B5\%D1\%81\%D1\%81\%D0\%B8\%D0\%B8_\%D0\%B7\%D0\%B0\%D0\%B4\%D0\%B0\%D1\%87_\%D0\%B8_\%D0\%BF\%D0\%BB\%D0\%B0\%D0\%BD\%D0\%B8\%D1\%80\%D0\%BE\%D0\%B2\%D1\%89\%D0\%B8\%D0\%BA.pdf\#nameddest=swr-ses-020-\%D0\%BF\%D1\%80\%D0\%B5\%D0\%B4\%D0\%B5\%D0\%BB-\%D0\%BE\%D0\%B1\%D1\%80\%D0\%B0\%D0\%B1\%D0\%BE\%D1\%82\%D0\%BA\%D0\%B8}{программном правиле предела обработки}.

Основание: правило ограниченной обработки в карточке домена; обработка технического сбоя по SR-SES-010.

Проверка: задержка проверяющего или очереди не оставляет задачу в бессрочном ожидании. После технического отказа повтор допустим до общего срока; после него задача завершается по просрочке. Вердикт, поступивший после прекращения обработки, не изменяет итог.

\subsection{SR-SES-013. Восстановление принятого ответа}\label{sr-ses-013-ux432ux43eux441ux441ux442ux430ux43dux43eux432ux43bux435ux43dux438ux435-ux43fux440ux438ux43dux44fux442ux43eux433ux43e-ux43eux442ux432ux435ux442ux430}

\hypertarget{sr-ses-013-восстановление-принятого-ответа}{}

После приёма своевременного ответа его дальнейшая судьба не должна зависеть только от памяти работающего процесса. До математического вычисления сохраняются входные данные и временные отметки. При перезапуске система различает операцию с уже записанным результатом и принятую операцию, которая ещё не закончилась.

\textbf{Требование.} После перезапуска система должна продолжать обработку сохранённого своевременного ответа, если его 60-секундный срок ещё не истёк, используя прежнее условие, ответ и время получения.

Операция без результата продолжается по тому же ответу и условию, если её собственный срок ещё действует. Истёкшая операция завершается техническим отказом; время поступления не заменяется временем запуска приложения. Уже записанный вердикт повторно не расходует попытку. Программная запись операции определена в \href{../07_\%D0\%9F\%D1\%80\%D0\%BE\%D0\%B3\%D1\%80\%D0\%B0\%D0\%BC\%D0\%BC\%D0\%BD\%D1\%8B\%D0\%B5_\%D1\%82\%D1\%80\%D0\%B5\%D0\%B1\%D0\%BE\%D0\%B2\%D0\%B0\%D0\%BD\%D0\%B8\%D1\%8F/DOC-06.2_\%D0\%A1\%D0\%B5\%D1\%81\%D1\%81\%D0\%B8\%D0\%B8_\%D0\%B7\%D0\%B0\%D0\%B4\%D0\%B0\%D1\%87_\%D0\%B8_\%D0\%BF\%D0\%BB\%D0\%B0\%D0\%BD\%D0\%B8\%D1\%80\%D0\%BE\%D0\%B2\%D1\%89\%D0\%B8\%D0\%BA.pdf\#nameddest=swr-ses-021-\%D0\%B7\%D0\%B0\%D0\%BF\%D0\%B8\%D1\%81\%D1\%8C-\%D0\%BF\%D1\%80\%D0\%B8\%D0\%BD\%D1\%8F\%D1\%82\%D0\%BE\%D0\%B9-\%D0\%BE\%D0\%BF\%D0\%B5\%D1\%80\%D0\%B0\%D1\%86\%D0\%B8\%D0\%B8}{требовании сохранения принятого ответа}, а получение готового итога --- в \hyperref[sr-ses-014-ux43fux43eux43bux443ux447ux435ux43dux438ux435-ux441ux43eux445ux440ux430ux43dux451ux43dux43dux43eux433ux43e-ux440ux435ux437ux443ux43bux44cux442ux430ux442ux430]{чтении результата}.

Основание: правило восстановления в карточке домена; \texttt{[TS.\allowbreak{}REQ.\allowbreak{}REL.\allowbreak{}PERSIST]}.

Проверка: перезапуск до фиксации результата не начинает заново срок сессии и обработки. Уже сохранённый вердикт не вычисляется повторно как новая попытка. Если срок обработки истёк, применяется SR-SES-012.

\subsection{SR-SES-014. Получение сохранённого результата}\label{sr-ses-014-ux43fux43eux43bux443ux447ux435ux43dux438ux435-ux441ux43eux445ux440ux430ux43dux451ux43dux43dux43eux433ux43e-ux440ux435ux437ux443ux43bux44cux442ux430ux442ux430}

\hypertarget{sr-ses-014-получение-сохранённого-результата}{}

Сохранение результата и доставка HTTP-ответа --- разные события. Соединение может пропасть после записи вердикта, и клиент останется без ответа, хотя задача уже завершена. Повторная отправка решения здесь непригодна: она должна быть отклонена как обращение к завершённой задаче. Для восстановления отображения используется отдельное чтение сессии.

\textbf{Требование.} Система должна предоставлять пользователю состояние собственной сессии и сохранённые результаты её задач по идентификатору сессии, в том числе после потери соединения или перезапуска.

Пользователь получает данные своего набора, состояния задач и их сохранённые предметные результаты. Срок ограничивает новые решения, но не закрывает чтение уже полученных итогов. Поэтому после перезапуска клиент может восстановить картину без математического пересчёта и новой попытки. Поля ответа описаны в \href{../07_\%D0\%9F\%D1\%80\%D0\%BE\%D0\%B3\%D1\%80\%D0\%B0\%D0\%BC\%D0\%BC\%D0\%BD\%D1\%8B\%D0\%B5_\%D1\%82\%D1\%80\%D0\%B5\%D0\%B1\%D0\%BE\%D0\%B2\%D0\%B0\%D0\%BD\%D0\%B8\%D1\%8F/DOC-06.2_\%D0\%A1\%D0\%B5\%D1\%81\%D1\%81\%D0\%B8\%D0\%B8_\%D0\%B7\%D0\%B0\%D0\%B4\%D0\%B0\%D1\%87_\%D0\%B8_\%D0\%BF\%D0\%BB\%D0\%B0\%D0\%BD\%D0\%B8\%D1\%80\%D0\%BE\%D0\%B2\%D1\%89\%D0\%B8\%D0\%BA.pdf\#nameddest=\%D1\%87\%D1\%82\%D0\%B5\%D0\%BD\%D0\%B8\%D0\%B5-\%D0\%BF\%D0\%BE\%D1\%81\%D0\%BB\%D0\%B5-\%D0\%BF\%D0\%BE\%D1\%82\%D0\%B5\%D1\%80\%D0\%B8-\%D0\%BE\%D1\%82\%D0\%B2\%D0\%B5\%D1\%82\%D0\%B0}{контракте чтения}, а программная операция --- в \href{../07_\%D0\%9F\%D1\%80\%D0\%BE\%D0\%B3\%D1\%80\%D0\%B0\%D0\%BC\%D0\%BC\%D0\%BD\%D1\%8B\%D0\%B5_\%D1\%82\%D1\%80\%D0\%B5\%D0\%B1\%D0\%BE\%D0\%B2\%D0\%B0\%D0\%BD\%D0\%B8\%D1\%8F/DOC-06.2_\%D0\%A1\%D0\%B5\%D1\%81\%D1\%81\%D0\%B8\%D0\%B8_\%D0\%B7\%D0\%B0\%D0\%B4\%D0\%B0\%D1\%87_\%D0\%B8_\%D0\%BF\%D0\%BB\%D0\%B0\%D0\%BD\%D0\%B8\%D1\%80\%D0\%BE\%D0\%B2\%D1\%89\%D0\%B8\%D0\%BA.pdf\#nameddest=swr-ses-023-\%D1\%87\%D1\%82\%D0\%B5\%D0\%BD\%D0\%B8\%D0\%B5-\%D1\%81\%D0\%B5\%D1\%81\%D1\%81\%D0\%B8\%D0\%B8-\%D0\%B8-\%D1\%80\%D0\%B5\%D0\%B7\%D1\%83\%D0\%BB\%D1\%8C\%D1\%82\%D0\%B0\%D1\%82\%D0\%B0}{требовании получения сессии}.

Основание: сохранность результатов по SR-SYS-007; сценарий восстановления карточки домена.

Проверка: вердикт сохранён, но ответ на запрос не доставлен клиенту. Чтение сессии возвращает тот же вердикт без новой предметной попытки. Получение результата после общего срока разрешено; чужая сессия недоступна.

\subsection{Пояснения к состояниям и попытке}\label{ux43fux43eux44fux441ux43dux435ux43dux438ux44f-ux43a-ux441ux43eux441ux442ux43eux44fux43dux438ux44fux43c-ux438-ux43fux43eux43fux44bux442ux43aux435}

\hypertarget{пояснения-к-состояниям-и-попытке}{}

Состояние \texttt{NEW} показывает отсутствие итогового результата. Оно не означает безусловного разрешения отправить ответ: дополнительно проверяются владелец, срок и наличие уже обрабатываемого ответа. Поэтому даже при задержке фонового завершения задача не получает дополнительного времени решения.

Предметная попытка завершается фиксацией результата проверки правильности. В неверно оформленном запросе такого результата нет; технический отказ также не устанавливает, правильно ли решена задача. Повтор после этих событий представляет новый запрос, своевременность которого проверяется заново. Он не наследует время поступления предыдущего запроса.

У своевременного допустимого ответа сохраняется право на завершение проверки. Состояние задачи в этот промежуток ещё не итоговое, но новая отправка не должна создавать вторую завершённую проверку. Внешнее обязательство системы состоит в сохранении единственного предметного результата; механизм согласования задаётся на программном уровне.

Просрочка и неверное решение приводят к одному итоговому состоянию \texttt{FAILED}, но основания различаются. Неверное решение связано с предметным вердиктом, а просрочка --- с прекращением возможности решения по общему сроку. Это различие сохраняется в описании проверок и диагностике операции.

\subsection{Сочетания событий для проверки раздела}\label{ux441ux43eux447ux435ux442ux430ux43dux438ux44f-ux441ux43eux431ux44bux442ux438ux439-ux434ux43bux44f-ux43fux440ux43eux432ux435ux440ux43aux438-ux440ux430ux437ux434ux435ux43bux430}

\hypertarget{сочетания-событий-для-проверки-раздела}{}

Приведённые сочетания раскрывают требования, а не вводят отдельные TC. При подготовке DOC-10 для них нужно задать исходное состояние, действие и наблюдаемый результат.

{\def\LTcaptype{none} % do not increment counter
\begin{longtable}[]{@{}
  >{\raggedright\arraybackslash}p{(\linewidth - 2\tabcolsep) * \real{0.3200}}
  >{\raggedright\arraybackslash}p{(\linewidth - 2\tabcolsep) * \real{0.6800}}@{}}
\toprule\noalign{}
\begin{minipage}[b]{\linewidth}\raggedright
Последовательность событий
\end{minipage} & \begin{minipage}[b]{\linewidth}\raggedright
Ожидаемый результат
\end{minipage} \\
\midrule\noalign{}
\endhead
\bottomrule\noalign{}
\endlastfoot
Создан набор; один ответ неверный; остальные задачи ещё \texttt{NEW} & Задача завершена с \texttt{FAILED}, сессия остаётся \texttt{NEW} \\
Все задачи проверены верно & Задачи и сессия имеют \texttt{SOLVED} \\
Все задачи завершены, хотя бы одна неверно & Итог сессии \texttt{FAILED} \\
Некорректный запрос получен до срока; исправленный запрос тоже получен до срока & Первый запрос не расходует попытку; второй допускается при соблюдении остальных условий \\
Технический сбой произошёл до срока; новый запрос получен после срока & Новый ответ отклоняется, несмотря на сохранённую предметную попытку \\
Допустимый ответ получен до срока; пока он проверяется, наступает срок & Задача ожидает вердикт, планировщик не завершает её по просрочке \\
Своевременная проверка после срока закончилась техническим сбоем & Предметного вердикта нет; новый ответ не допускается, задача завершается по просрочке \\
Отправлен ответ на задачу с итоговым состоянием & Результат сохраняется без повторной предметной проверки \\
\end{longtable}
}

Для общего срока важно рассматривать не только момент запроса, но и состояние других задач сессии. Завершение одной проверки не должно закрывать сессию, пока для другой задачи остаётся своевременный ответ в обработке. Аналогично окончание срока не изменяет уже определённые результаты.

Редакция 07.11.2026: развёрнуты пояснения каждого требования, добавлены примеры и точные ссылки на связанные разделы. Формулировки обязательств и идентификаторы сохранены.

\section{DOC-05. Требования к интеграции}\label{doc-05-ux442ux440ux435ux431ux43eux432ux430ux43dux438ux44f-ux43a-ux438ux43dux442ux435ux433ux440ux430ux446ux438ux438}

\hypertarget{doc-05-требования-к-интеграции}{}

Раздел \hyperref[doc-05-ux441ux438ux441ux442ux435ux43cux43dux44bux435-ux442ux440ux435ux431ux43eux432ux430ux43dux438ux44f]{DOC-05}, версия 0.4 от 07.11.2026. Автор: Анохин Г.А. Приоритет всех требований --- обязательный. Здесь определены свойства совместной работы частей системы. Имена полей, HTTP-коды и программные механизмы описываются ниже по иерархии.

\subsection{SR-SYS-001. Единая личность пользователя}\label{sr-sys-001-ux435ux434ux438ux43dux430ux44f-ux43bux438ux447ux43dux43eux441ux442ux44c-ux43fux43eux43bux44cux437ux43eux432ux430ux442ux435ux43bux44f}

\hypertarget{sr-sys-001-единая-личность-пользователя}{}

Сессии, группы и конфигурации используют одну личность пользователя, установленную проверкой доступа. Иначе одна часть могла бы выбирать конфигурацию для одного человека, а другая сохранять результат от имени другого. Пользовательские параметры определяют объект операции, но не меняют её исполнителя. Передачу этих сведений между частями показывает \href{../07_\%D0\%9F\%D1\%80\%D0\%BE\%D0\%B3\%D1\%80\%D0\%B0\%D0\%BC\%D0\%BC\%D0\%BD\%D1\%8B\%D0\%B5_\%D1\%82\%D1\%80\%D0\%B5\%D0\%B1\%D0\%BE\%D0\%B2\%D0\%B0\%D0\%BD\%D0\%B8\%D1\%8F/DOC-06_\%D0\%9F\%D1\%80\%D0\%BE\%D0\%B3\%D1\%80\%D0\%B0\%D0\%BC\%D0\%BC\%D0\%BD\%D1\%8B\%D0\%B5_\%D1\%82\%D1\%80\%D0\%B5\%D0\%B1\%D0\%BE\%D0\%B2\%D0\%B0\%D0\%BD\%D0\%B8\%D1\%8F.pdf\#nameddest=2-\%D0\%BE\%D1\%82\%D0\%B2\%D0\%B5\%D1\%82\%D1\%81\%D1\%82\%D0\%B2\%D0\%B5\%D0\%BD\%D0\%BD\%D0\%BE\%D1\%81\%D1\%82\%D1\%8C-\%D0\%B8-\%D1\%81\%D0\%B2\%D1\%8F\%D0\%B7\%D0\%B8}{описание ответственности компонентов}.

\textbf{Требование.} Система должна использовать подтверждённую личность текущего пользователя при проверке принадлежности сессий и доступности конфигураций.

При создании набора по подтверждённому пользователю определяется его группа. При ответе тот же пользователь сопоставляется с владельцем сохранённой сессии. Подмена \texttt{userId} в теле не должна менять эти решения. Это системное обязательство раскрывается в \href{../07_\%D0\%9F\%D1\%80\%D0\%BE\%D0\%B3\%D1\%80\%D0\%B0\%D0\%BC\%D0\%BC\%D0\%BD\%D1\%8B\%D0\%B5_\%D1\%82\%D1\%80\%D0\%B5\%D0\%B1\%D0\%BE\%D0\%B2\%D0\%B0\%D0\%BD\%D0\%B8\%D1\%8F/DOC-06_\%D0\%9F\%D1\%80\%D0\%BE\%D0\%B3\%D1\%80\%D0\%B0\%D0\%BC\%D0\%BC\%D0\%BD\%D1\%8B\%D0\%B5_\%D1\%82\%D1\%80\%D0\%B5\%D0\%B1\%D0\%BE\%D0\%B2\%D0\%B0\%D0\%BD\%D0\%B8\%D1\%8F.pdf\#nameddest=swr-sys-001-\%D0\%BF\%D0\%B5\%D1\%80\%D0\%B5\%D0\%B4\%D0\%B0\%D1\%87\%D0\%B0-\%D0\%BB\%D0\%B8\%D1\%87\%D0\%BD\%D0\%BE\%D1\%81\%D1\%82\%D0\%B8-\%D0\%BF\%D0\%BE\%D0\%BB\%D1\%8C\%D0\%B7\%D0\%BE\%D0\%B2\%D0\%B0\%D1\%82\%D0\%B5\%D0\%BB\%D1\%8F}{общем правиле передачи личности} и в \href{../07_\%D0\%9F\%D1\%80\%D0\%BE\%D0\%B3\%D1\%80\%D0\%B0\%D0\%BC\%D0\%BC\%D0\%BD\%D1\%8B\%D0\%B5_\%D1\%82\%D1\%80\%D0\%B5\%D0\%B1\%D0\%BE\%D0\%B2\%D0\%B0\%D0\%BD\%D0\%B8\%D1\%8F/DOC-06.2_\%D0\%A1\%D0\%B5\%D1\%81\%D1\%81\%D0\%B8\%D0\%B8_\%D0\%B7\%D0\%B0\%D0\%B4\%D0\%B0\%D1\%87_\%D0\%B8_\%D0\%BF\%D0\%BB\%D0\%B0\%D0\%BD\%D0\%B8\%D1\%80\%D0\%BE\%D0\%B2\%D1\%89\%D0\%B8\%D0\%BA.pdf\#nameddest=swr-ses-001-\%D0\%BF\%D0\%BE\%D0\%BB\%D1\%83\%D1\%87\%D0\%B5\%D0\%BD\%D0\%B8\%D0\%B5-\%D0\%BF\%D0\%BE\%D0\%BB\%D1\%8C\%D0\%B7\%D0\%BE\%D0\%B2\%D0\%B0\%D1\%82\%D0\%B5\%D0\%BB\%D1\%8F}{получении пользователя обработчиком сессий}.

Основание: \texttt{[TS.\allowbreak{}REQ.\allowbreak{}FUNC.\allowbreak{}AUTH.\allowbreak{}JWT]}, \texttt{[TS.\allowbreak{}REQ.\allowbreak{}FUNC.\allowbreak{}SESSION.\allowbreak{}SUBMIT]} в \href{../02_\%D0\%A2\%D0\%B5\%D1\%85\%D0\%BD\%D0\%B8\%D1\%87\%D0\%B5\%D1\%81\%D0\%BA\%D0\%BE\%D0\%B5_\%D0\%B7\%D0\%B0\%D0\%B4\%D0\%B0\%D0\%BD\%D0\%B8\%D0\%B5/\%D0\%A2\%D0\%B5\%D1\%85\%D0\%BD\%D0\%B8\%D1\%87\%D0\%B5\%D1\%81\%D0\%BA\%D0\%BE\%D0\%B5_\%D0\%B7\%D0\%B0\%D0\%B4\%D0\%B0\%D0\%BD\%D0\%B8\%D0\%B5.pdf}{ТЗ}.

Проверка: указание чужого пользователя в параметрах запроса не даёт доступа к его сессиям и конфигурациям. Данные аутентификации и владельца сессии сопоставляются однозначно.

\subsection{SR-SYS-002. Связность данных задачи}\label{sr-sys-002-ux441ux432ux44fux437ux43dux43eux441ux442ux44c-ux434ux430ux43dux43dux44bux445-ux437ux430ux434ux430ux447ux438}

\hypertarget{sr-sys-002-связность-данных-задачи}{}

Идентификаторы связывают операции, выполняемые в разное время. Пользователь получает задачу при создании набора, отправляет её ID при ответе и позднее читает результат. На всех этапах речь должна идти об одном объекте. Сохранённая принадлежность задачи сессии одновременно используется для поиска результата и проверки доступа.

\textbf{Требование.} Система должна сохранять однозначные связи задачи с сессией и сессии с пользователем при создании набора, проверке ответа и получении результата.

Тип задачи или текст условия не заменяет её ID: в одном наборе могут быть задания одного типа, а данные разных пользователей могут иметь сходный вид. Связь устанавливается сохранёнными идентификаторами. Передаваемые поля разобраны в \href{../07_\%D0\%9F\%D1\%80\%D0\%BE\%D0\%B3\%D1\%80\%D0\%B0\%D0\%BC\%D0\%BC\%D0\%BD\%D1\%8B\%D0\%B5_\%D1\%82\%D1\%80\%D0\%B5\%D0\%B1\%D0\%BE\%D0\%B2\%D0\%B0\%D0\%BD\%D0\%B8\%D1\%8F/DOC-06.2_\%D0\%A1\%D0\%B5\%D1\%81\%D1\%81\%D0\%B8\%D0\%B8_\%D0\%B7\%D0\%B0\%D0\%B4\%D0\%B0\%D1\%87_\%D0\%B8_\%D0\%BF\%D0\%BB\%D0\%B0\%D0\%BD\%D0\%B8\%D1\%80\%D0\%BE\%D0\%B2\%D1\%89\%D0\%B8\%D0\%BA.pdf\#nameddest=\%D1\%80\%D0\%B5\%D0\%B7\%D1\%83\%D0\%BB\%D1\%8C\%D1\%82\%D0\%B0\%D1\%82-\%D0\%B8-\%D0\%BE\%D1\%82\%D0\%BA\%D0\%B0\%D0\%B7\%D1\%8B}{контракте результата}, а единое использование идентификаторов компонентами закреплено в \href{../07_\%D0\%9F\%D1\%80\%D0\%BE\%D0\%B3\%D1\%80\%D0\%B0\%D0\%BC\%D0\%BC\%D0\%BD\%D1\%8B\%D0\%B5_\%D1\%82\%D1\%80\%D0\%B5\%D0\%B1\%D0\%BE\%D0\%B2\%D0\%B0\%D0\%BD\%D0\%B8\%D1\%8F/DOC-06_\%D0\%9F\%D1\%80\%D0\%BE\%D0\%B3\%D1\%80\%D0\%B0\%D0\%BC\%D0\%BC\%D0\%BD\%D1\%8B\%D0\%B5_\%D1\%82\%D1\%80\%D0\%B5\%D0\%B1\%D0\%BE\%D0\%B2\%D0\%B0\%D0\%BD\%D0\%B8\%D1\%8F.pdf\#nameddest=swr-sys-002-\%D0\%BE\%D0\%B1\%D1\%89\%D0\%B8\%D0\%B5-\%D0\%B8\%D0\%B4\%D0\%B5\%D0\%BD\%D1\%82\%D0\%B8\%D1\%84\%D0\%B8\%D0\%BA\%D0\%B0\%D1\%82\%D0\%BE\%D1\%80\%D1\%8B}{программном требовании общих ID}.

Основание: \texttt{[TS.\allowbreak{}REQ.\allowbreak{}FUNC.\allowbreak{}SESSION.\allowbreak{}CREATE]}, \texttt{[TS.\allowbreak{}REQ.\allowbreak{}FUNC.\allowbreak{}SESSION.\allowbreak{}SUBMIT]}.

Проверка: задача в результате создания и в результате её проверки относится к одной сессии и владельцу; операция над одной задачей не изменяет данные другой.

\subsection{SR-SYS-003. Согласованный состав набора}\label{sr-sys-003-ux441ux43eux433ux43bux430ux441ux43eux432ux430ux43dux43dux44bux439-ux441ux43eux441ux442ux430ux432-ux43dux430ux431ux43eux440ux430}

\hypertarget{sr-sys-003-согласованный-состав-набора}{}

Конфигурация описывает состав нового набора и его доступность группе. Сессия фиксирует конкретные задания, которые уже выданы пользователю. Эти объекты связаны в момент создания, но дальнейшая проверка не должна заново определять условие по изменившимся настройкам. Граница между выбором состава и жизненным циклом набора проведена в \hyperref[sr-ses-001-ux441ux43eux437ux434ux430ux43dux438ux435-ux43dux430ux431ux43eux440ux430-ux43fux43e-ux43aux43eux43dux444ux438ux433ux443ux440ux430ux446ux438ux438-ux433ux440ux443ux43fux43fux44b]{правиле создания сессии}.

\textbf{Требование.} Система должна формировать состав новой сессии по выбранной конфигурации, доступной группе пользователя на момент создания.

Если преподаватель изменил доступный состав после выдачи задач, проверяется прежнее выданное условие. Иначе студент мог бы получить вердикт по другому заданию. Система сохраняет состав и условия при создании, а недоступную конфигурацию отклоняет до подтверждения операции. Программная детализация находится в \href{../07_\%D0\%9F\%D1\%80\%D0\%BE\%D0\%B3\%D1\%80\%D0\%B0\%D0\%BC\%D0\%BC\%D0\%BD\%D1\%8B\%D0\%B5_\%D1\%82\%D1\%80\%D0\%B5\%D0\%B1\%D0\%BE\%D0\%B2\%D0\%B0\%D0\%BD\%D0\%B8\%D1\%8F/DOC-06_\%D0\%9F\%D1\%80\%D0\%BE\%D0\%B3\%D1\%80\%D0\%B0\%D0\%BC\%D0\%BC\%D0\%BD\%D1\%8B\%D0\%B5_\%D1\%82\%D1\%80\%D0\%B5\%D0\%B1\%D0\%BE\%D0\%B2\%D0\%B0\%D0\%BD\%D0\%B8\%D1\%8F.pdf\#nameddest=swr-sys-003-\%D1\%81\%D0\%BE\%D1\%81\%D1\%82\%D0\%B0\%D0\%B2-\%D1\%81\%D0\%BE\%D0\%B7\%D0\%B4\%D0\%B0\%D0\%BD\%D0\%BD\%D0\%BE\%D0\%B9-\%D1\%81\%D0\%B5\%D1\%81\%D1\%81\%D0\%B8\%D0\%B8}{сохранении состава} и \href{../07_\%D0\%9F\%D1\%80\%D0\%BE\%D0\%B3\%D1\%80\%D0\%B0\%D0\%BC\%D0\%BC\%D0\%BD\%D1\%8B\%D0\%B5_\%D1\%82\%D1\%80\%D0\%B5\%D0\%B1\%D0\%BE\%D0\%B2\%D0\%B0\%D0\%BD\%D0\%B8\%D1\%8F/DOC-06.2_\%D0\%A1\%D0\%B5\%D1\%81\%D1\%81\%D0\%B8\%D0\%B8_\%D0\%B7\%D0\%B0\%D0\%B4\%D0\%B0\%D1\%87_\%D0\%B8_\%D0\%BF\%D0\%BB\%D0\%B0\%D0\%BD\%D0\%B8\%D1\%80\%D0\%BE\%D0\%B2\%D1\%89\%D0\%B8\%D0\%BA.pdf\#nameddest=swr-ses-004-\%D1\%84\%D0\%BE\%D1\%80\%D0\%BC\%D0\%B8\%D1\%80\%D0\%BE\%D0\%B2\%D0\%B0\%D0\%BD\%D0\%B8\%D0\%B5-\%D0\%B7\%D0\%B0\%D0\%B4\%D0\%B0\%D1\%87}{формировании задач}.

Основание: \texttt{[TS.\allowbreak{}REQ.\allowbreak{}FUNC.\allowbreak{}CONFIG.\allowbreak{}VIEW]}, \texttt{[TS.\allowbreak{}REQ.\allowbreak{}FUNC.\allowbreak{}SESSION.\allowbreak{}CREATE]}.

Проверка: состав соответствует включённой конфигурации. Конфигурация другой группы или выключенная конфигурация не используется. Последующее изменение настроек не подменяет условия уже выданных задач.

\subsection{SR-SYS-004. Соответствие генерации и проверки}\label{sr-sys-004-ux441ux43eux43eux442ux432ux435ux442ux441ux442ux432ux438ux435-ux433ux435ux43dux435ux440ux430ux446ux438ux438-ux438-ux43fux440ux43eux432ux435ux440ux43aux438}

\hypertarget{sr-sys-004-соответствие-генерации-и-проверки}{}

Генерация и проверка принадлежат одному предметному контракту. Генератор формирует условие определённого типа, а проверяющий должен интерпретировать именно это условие вместе с ответом. Домен сессий соединяет эти операции, сохраняя тип и условие задачи между выдачей и отправкой решения.

\textbf{Требование.} Система должна проверять ответ по тому же типу задачи и сохранённому условию, которые использовались при её выдаче.

Вызвать существующий проверяющий недостаточно: он должен соответствовать выбранной задаче. Клиент не может переключить этот выбор произвольным значением поля \texttt{type}. Если обработчик нужного типа недоступен, результат относится к техническому отказу. Форматы определены в \href{../07_\%D0\%9F\%D1\%80\%D0\%BE\%D0\%B3\%D1\%80\%D0\%B0\%D0\%BC\%D0\%BC\%D0\%BD\%D1\%8B\%D0\%B5_\%D1\%82\%D1\%80\%D0\%B5\%D0\%B1\%D0\%BE\%D0\%B2\%D0\%B0\%D0\%BD\%D0\%B8\%D1\%8F/DOC-06.2_\%D0\%A1\%D0\%B5\%D1\%81\%D1\%81\%D0\%B8\%D0\%B8_\%D0\%B7\%D0\%B0\%D0\%B4\%D0\%B0\%D1\%87_\%D0\%B8_\%D0\%BF\%D0\%BB\%D0\%B0\%D0\%BD\%D0\%B8\%D1\%80\%D0\%BE\%D0\%B2\%D1\%89\%D0\%B8\%D0\%BA.pdf\#nameddest=\%D0\%BF\%D1\%80\%D0\%B5\%D0\%B4\%D0\%BC\%D0\%B5\%D1\%82\%D0\%BD\%D1\%8B\%D0\%B9-\%D0\%BE\%D1\%82\%D0\%B2\%D0\%B5\%D1\%82}{контракте предметного ответа}, передача данных --- в \href{../07_\%D0\%9F\%D1\%80\%D0\%BE\%D0\%B3\%D1\%80\%D0\%B0\%D0\%BC\%D0\%BC\%D0\%BD\%D1\%8B\%D0\%B5_\%D1\%82\%D1\%80\%D0\%B5\%D0\%B1\%D0\%BE\%D0\%B2\%D0\%B0\%D0\%BD\%D0\%B8\%D1\%8F/DOC-06.2_\%D0\%A1\%D0\%B5\%D1\%81\%D1\%81\%D0\%B8\%D0\%B8_\%D0\%B7\%D0\%B0\%D0\%B4\%D0\%B0\%D1\%87_\%D0\%B8_\%D0\%BF\%D0\%BB\%D0\%B0\%D0\%BD\%D0\%B8\%D1\%80\%D0\%BE\%D0\%B2\%D1\%89\%D0\%B8\%D0\%BA.pdf\#nameddest=swr-ses-010-\%D0\%BF\%D0\%B5\%D1\%80\%D0\%B5\%D0\%B4\%D0\%B0\%D1\%87\%D0\%B0-\%D0\%BF\%D1\%80\%D0\%BE\%D1\%84\%D0\%B8\%D0\%BB\%D1\%8C\%D0\%BD\%D0\%BE\%D0\%B9-\%D0\%BF\%D1\%80\%D0\%BE\%D0\%B2\%D0\%B5\%D1\%80\%D0\%BA\%D0\%B5}{правиле вызова профильной проверки}.

Основание: \texttt{[TS.\allowbreak{}REQ.\allowbreak{}FUNC.\allowbreak{}SESSION.\allowbreak{}CREATE]}, \texttt{[TS.\allowbreak{}REQ.\allowbreak{}FUNC.\allowbreak{}SESSION.\allowbreak{}SUBMIT]}; ответственность частей в DOC-05.

Проверка: решение одной задачи не проверяется с условием другой; обработчик неподдерживаемого типа не подменяется обработчиком другого типа.

\subsection{SR-SYS-005. Согласованность результата и состояния}\label{sr-sys-005-ux441ux43eux433ux43bux430ux441ux43eux432ux430ux43dux43dux43eux441ux442ux44c-ux440ux435ux437ux443ux43bux44cux442ux430ux442ux430-ux438-ux441ux43eux441ux442ux43eux44fux43dux438ux44f}

\hypertarget{sr-sys-005-согласованность-результата-и-состояния}{}

Вердикт имеет значение для пользователя только вместе с сохранённым состоянием. При успешной доставке ответа система подтверждает не просто выполненное вычисление, а завершённую операцию над задачей и сессией. Это связывает предметную проверку, хранение и внешний интерфейс. Раздельное рассмотрение этих частей не должно приводить к противоречащим итогам.

\textbf{Требование.} Система должна выдавать подтверждённый результат предметной проверки только при сохранении соответствующего состояния задачи и связанного итога сессии.

Если после неверного решения завершились все задачи, должны сохраниться и отрицательный результат задачи, и \texttt{FAILED} сессии. Если запись не завершилась, нельзя сообщать операцию как подтверждённую. После потери соединения уже сохранённый результат доступен через \hyperref[sr-ses-014-ux43fux43eux43bux443ux447ux435ux43dux438ux435-ux441ux43eux445ux440ux430ux43dux451ux43dux43dux43eux433ux43e-ux440ux435ux437ux443ux43bux44cux442ux430ux442ux430]{чтение сессии}. Границу записи определяют \href{../07_\%D0\%9F\%D1\%80\%D0\%BE\%D0\%B3\%D1\%80\%D0\%B0\%D0\%BC\%D0\%BC\%D0\%BD\%D1\%8B\%D0\%B5_\%D1\%82\%D1\%80\%D0\%B5\%D0\%B1\%D0\%BE\%D0\%B2\%D0\%B0\%D0\%BD\%D0\%B8\%D1\%8F/DOC-06_\%D0\%9F\%D1\%80\%D0\%BE\%D0\%B3\%D1\%80\%D0\%B0\%D0\%BC\%D0\%BC\%D0\%BD\%D1\%8B\%D0\%B5_\%D1\%82\%D1\%80\%D0\%B5\%D0\%B1\%D0\%BE\%D0\%B2\%D0\%B0\%D0\%BD\%D0\%B8\%D1\%8F.pdf\#nameddest=swr-sys-005-\%D1\%81\%D0\%BE\%D0\%B3\%D0\%BB\%D0\%B0\%D1\%81\%D0\%BE\%D0\%B2\%D0\%B0\%D0\%BD\%D0\%BD\%D0\%BE\%D0\%B5-\%D1\%81\%D0\%BE\%D1\%85\%D1\%80\%D0\%B0\%D0\%BD\%D0\%B5\%D0\%BD\%D0\%B8\%D0\%B5}{согласованное сохранение} и \href{../07_\%D0\%9F\%D1\%80\%D0\%BE\%D0\%B3\%D1\%80\%D0\%B0\%D0\%BC\%D0\%BC\%D0\%BD\%D1\%8B\%D0\%B5_\%D1\%82\%D1\%80\%D0\%B5\%D0\%B1\%D0\%BE\%D0\%B2\%D0\%B0\%D0\%BD\%D0\%B8\%D1\%8F/DOC-06.2_\%D0\%A1\%D0\%B5\%D1\%81\%D1\%81\%D0\%B8\%D0\%B8_\%D0\%B7\%D0\%B0\%D0\%B4\%D0\%B0\%D1\%87_\%D0\%B8_\%D0\%BF\%D0\%BB\%D0\%B0\%D0\%BD\%D0\%B8\%D1\%80\%D0\%BE\%D0\%B2\%D1\%89\%D0\%B8\%D0\%BA.pdf\#nameddest=swr-ses-011-\%D1\%81\%D0\%BE\%D1\%85\%D1\%80\%D0\%B0\%D0\%BD\%D0\%B5\%D0\%BD\%D0\%B8\%D0\%B5-\%D0\%B2\%D0\%B5\%D1\%80\%D0\%B4\%D0\%B8\%D0\%BA\%D1\%82\%D0\%B0}{фиксация вердикта}.

Основание: \texttt{[TS.\allowbreak{}REQ.\allowbreak{}FUNC.\allowbreak{}SESSION.\allowbreak{}SUBMIT]}, \texttt{[TS.\allowbreak{}REQ.\allowbreak{}REL.\allowbreak{}PERSIST]}.

Проверка: ошибка сохранения не возвращается как успешное завершение операции; после завершённой операции выдаваемый результат и сохранённые состояния согласованы.

\subsection{SR-SYS-006. Единое правило времени}\label{sr-sys-006-ux435ux434ux438ux43dux43eux435-ux43fux440ux430ux432ux438ux43bux43e-ux432ux440ux435ux43cux435ux43dux438}

\hypertarget{sr-sys-006-единое-правило-времени}{}

Обработка ответа и планировщик наблюдают одни задачи, но запускаются по разным причинам. Обработчик реагирует на запрос пользователя, планировщик --- на наступление времени. Для одной задачи их решения должны использовать один срок и одинаково учитывать своевременно полученный ответ. Собственная независимая трактовка срока у каждого компонента здесь недопустима.

\textbf{Требование.} Система должна применять общий серверный срок сессии при обработке ответов и при фоновом завершении задач.

Фоновая задержка не разрешает поздний запрос, а наступление общего срока не отменяет уже принятую своевременную проверку. Этот общий результат важнее порядка фактического запуска компонентов. Временная модель дана в \href{../notes/domains/\%D0\%A1\%D0\%B5\%D1\%81\%D1\%81\%D0\%B8\%D0\%B8_\%D0\%B7\%D0\%B0\%D0\%B4\%D0\%B0\%D1\%87.pdf\#nameddest=\%D1\%81\%D1\%80\%D0\%BE\%D0\%BA-\%D0\%B8-\%D0\%BF\%D0\%BE\%D0\%BF\%D1\%8B\%D1\%82\%D0\%BA\%D0\%B0}{карточке домена}, условия завершения --- в \href{../07_\%D0\%9F\%D1\%80\%D0\%BE\%D0\%B3\%D1\%80\%D0\%B0\%D0\%BC\%D0\%BC\%D0\%BD\%D1\%8B\%D0\%B5_\%D1\%82\%D1\%80\%D0\%B5\%D0\%B1\%D0\%BE\%D0\%B2\%D0\%B0\%D0\%BD\%D0\%B8\%D1\%8F/DOC-06.2_\%D0\%A1\%D0\%B5\%D1\%81\%D1\%81\%D0\%B8\%D0\%B8_\%D0\%B7\%D0\%B0\%D0\%B4\%D0\%B0\%D1\%87_\%D0\%B8_\%D0\%BF\%D0\%BB\%D0\%B0\%D0\%BD\%D0\%B8\%D1\%80\%D0\%BE\%D0\%B2\%D1\%89\%D0\%B8\%D0\%BA.pdf\#nameddest=swr-ses-017-\%D0\%B7\%D0\%B0\%D0\%B2\%D0\%B5\%D1\%80\%D1\%88\%D0\%B5\%D0\%BD\%D0\%B8\%D0\%B5-\%D0\%BF\%D0\%BE-\%D0\%BF\%D1\%80\%D0\%BE\%D1\%81\%D1\%80\%D0\%BE\%D1\%87\%D0\%BA\%D0\%B5}{обработке просрочки}, а защита своевременной операции --- в \href{../07_\%D0\%9F\%D1\%80\%D0\%BE\%D0\%B3\%D1\%80\%D0\%B0\%D0\%BC\%D0\%BC\%D0\%BD\%D1\%8B\%D0\%B5_\%D1\%82\%D1\%80\%D0\%B5\%D0\%B1\%D0\%BE\%D0\%B2\%D0\%B0\%D0\%BD\%D0\%B8\%D1\%8F/DOC-06.2_\%D0\%A1\%D0\%B5\%D1\%81\%D1\%81\%D0\%B8\%D0\%B8_\%D0\%B7\%D0\%B0\%D0\%B4\%D0\%B0\%D1\%87_\%D0\%B8_\%D0\%BF\%D0\%BB\%D0\%B0\%D0\%BD\%D0\%B8\%D1\%80\%D0\%BE\%D0\%B2\%D1\%89\%D0\%B8\%D0\%BA.pdf\#nameddest=swr-ses-015-\%D0\%B7\%D0\%B0\%D0\%B2\%D0\%B5\%D1\%80\%D1\%88\%D0\%B5\%D0\%BD\%D0\%B8\%D0\%B5-\%D1\%81\%D0\%B2\%D0\%BE\%D0\%B5\%D0\%B2\%D1\%80\%D0\%B5\%D0\%BC\%D0\%B5\%D0\%BD\%D0\%BD\%D0\%BE\%D0\%B9-\%D0\%BF\%D1\%80\%D0\%BE\%D0\%B2\%D0\%B5\%D1\%80\%D0\%BA\%D0\%B8}{правиле её завершения}.

Основание: \texttt{{[}TS.REQ.REL.SCHED{]}}; уточнённое правило срока из \href{../notes/domains/\%D0\%A1\%D0\%B5\%D1\%81\%D1\%81\%D0\%B8\%D0\%B8_\%D0\%B7\%D0\%B0\%D0\%B4\%D0\%B0\%D1\%87.pdf\#nameddest=\%D0\%BA\%D0\%B0\%D1\%80\%D1\%82\%D0\%BE\%D1\%87\%D0\%BA\%D0\%B0-\%D0\%B4\%D0\%BE\%D0\%BC\%D0\%B5\%D0\%BD\%D0\%B0-\%D1\%81\%D0\%B5\%D1\%81\%D1\%81\%D0\%B8\%D0\%B8-\%D0\%B7\%D0\%B0\%D0\%B4\%D0\%B0\%D1\%87}{карточки домена}.

Проверка: задержка планировщика не разрешает поздний ответ. Фоновая обработка не отменяет вердикт своевременно полученного ответа.

\subsection{SR-SYS-007. Сохранность данных сессии}\label{sr-sys-007-ux441ux43eux445ux440ux430ux43dux43dux43eux441ux442ux44c-ux434ux430ux43dux43dux44bux445-ux441ux435ux441ux441ux438ux438}

\hypertarget{sr-sys-007-сохранность-данных-сессии}{}

Время создания, принадлежность и результаты являются данными сессии, а не временными сведениями интерфейса. Потеря этих данных при перезапуске меняла бы условия решения и делала невозможным получение прежнего итога. Поэтому сохранность охватывает не только состояние \texttt{SOLVED} или \texttt{FAILED}, но и сведения, по которым это состояние относится к конкретному набору.

\textbf{Требование.} Система должна сохранять принадлежность, состав, время создания и определённые результаты задач при перезапуске приложения.

Завершённый результат читается без повторного вычисления. Для принятого ответа без результата применяется другой путь: восстановление по сохранённой записи с прежним сроком обработки. Эти случаи разделены в \hyperref[sr-ses-013-ux432ux43eux441ux441ux442ux430ux43dux43eux432ux43bux435ux43dux438ux435-ux43fux440ux438ux43dux44fux442ux43eux433ux43e-ux43eux442ux432ux435ux442ux430]{восстановлении принятого ответа} и \hyperref[sr-ses-014-ux43fux43eux43bux443ux447ux435ux43dux438ux435-ux441ux43eux445ux440ux430ux43dux451ux43dux43dux43eux433ux43e-ux440ux435ux437ux443ux43bux44cux442ux430ux442ux430]{получении готового результата}. Программное хранение описывают \href{../07_\%D0\%9F\%D1\%80\%D0\%BE\%D0\%B3\%D1\%80\%D0\%B0\%D0\%BC\%D0\%BC\%D0\%BD\%D1\%8B\%D0\%B5_\%D1\%82\%D1\%80\%D0\%B5\%D0\%B1\%D0\%BE\%D0\%B2\%D0\%B0\%D0\%BD\%D0\%B8\%D1\%8F/DOC-06_\%D0\%9F\%D1\%80\%D0\%BE\%D0\%B3\%D1\%80\%D0\%B0\%D0\%BC\%D0\%BC\%D0\%BD\%D1\%8B\%D0\%B5_\%D1\%82\%D1\%80\%D0\%B5\%D0\%B1\%D0\%BE\%D0\%B2\%D0\%B0\%D0\%BD\%D0\%B8\%D1\%8F.pdf\#nameddest=swr-sys-008-\%D1\%81\%D0\%BE\%D1\%85\%D1\%80\%D0\%B0\%D0\%BD\%D0\%BD\%D0\%BE\%D1\%81\%D1\%82\%D1\%8C-\%D0\%B7\%D0\%B0\%D0\%B2\%D0\%B5\%D1\%80\%D1\%88\%D1\%91\%D0\%BD\%D0\%BD\%D1\%8B\%D1\%85-\%D0\%B4\%D0\%B0\%D0\%BD\%D0\%BD\%D1\%8B\%D1\%85}{общая сохранность} и \href{../07_\%D0\%9F\%D1\%80\%D0\%BE\%D0\%B3\%D1\%80\%D0\%B0\%D0\%BC\%D0\%BC\%D0\%BD\%D1\%8B\%D0\%B5_\%D1\%82\%D1\%80\%D0\%B5\%D0\%B1\%D0\%BE\%D0\%B2\%D0\%B0\%D0\%BD\%D0\%B8\%D1\%8F/DOC-06.2_\%D0\%A1\%D0\%B5\%D1\%81\%D1\%81\%D0\%B8\%D0\%B8_\%D0\%B7\%D0\%B0\%D0\%B4\%D0\%B0\%D1\%87_\%D0\%B8_\%D0\%BF\%D0\%BB\%D0\%B0\%D0\%BD\%D0\%B8\%D1\%80\%D0\%BE\%D0\%B2\%D1\%89\%D0\%B8\%D0\%BA.pdf\#nameddest=swr-ses-021-\%D0\%B7\%D0\%B0\%D0\%BF\%D0\%B8\%D1\%81\%D1\%8C-\%D0\%BF\%D1\%80\%D0\%B8\%D0\%BD\%D1\%8F\%D1\%82\%D0\%BE\%D0\%B9-\%D0\%BE\%D0\%BF\%D0\%B5\%D1\%80\%D0\%B0\%D1\%86\%D0\%B8\%D0\%B8}{запись операции}.

Основание: \texttt{[TS.\allowbreak{}REQ.\allowbreak{}REL.\allowbreak{}PERSIST]}.

Проверка: после перезапуска срок не начинается заново, завершённые результаты доступны и не пересчитываются. Незавершённый сохранённый ответ обрабатывается по SR-SES-013 с учётом прежнего срока обработки.

\subsection{SR-SYS-008. Совместимость обмена}\label{sr-sys-008-ux441ux43eux432ux43cux435ux441ux442ux438ux43cux43eux441ux442ux44c-ux43eux431ux43cux435ux43dux430}

\hypertarget{sr-sys-008-совместимость-обмена}{}

Части системы должны одинаково понимать данные на внешней границе. При создании клиент получает ID, тип и условие задачи; при отправке использует их для выбора задачи и схемы ответа; после проверки сопоставляет результат с тем же заданием. Совместимость означает сохранение этого смысла на всём пути, а не только возможность разобрать отдельный JSON.

\textbf{Требование.} Система должна обмениваться данными пользовательских операций через REST-интерфейс с согласованными JSON-форматами запросов и результатов.

Поле \texttt{statementId} ссылается на ID задачи из выданного набора. Большие целые передаются без потери точности, а время однозначно задаёт серверный момент. Для работы клиента эти детали собраны в \href{../07_\%D0\%9F\%D1\%80\%D0\%BE\%D0\%B3\%D1\%80\%D0\%B0\%D0\%BC\%D0\%BC\%D0\%BD\%D1\%8B\%D0\%B5_\%D1\%82\%D1\%80\%D0\%B5\%D0\%B1\%D0\%BE\%D0\%B2\%D0\%B0\%D0\%BD\%D0\%B8\%D1\%8F/DOC-06.2_\%D0\%A1\%D0\%B5\%D1\%81\%D1\%81\%D0\%B8\%D0\%B8_\%D0\%B7\%D0\%B0\%D0\%B4\%D0\%B0\%D1\%87_\%D0\%B8_\%D0\%BF\%D0\%BB\%D0\%B0\%D0\%BD\%D0\%B8\%D1\%80\%D0\%BE\%D0\%B2\%D1\%89\%D0\%B8\%D0\%BA.pdf\#nameddest=\%D0\%BE\%D0\%BF\%D0\%B5\%D1\%80\%D0\%B0\%D1\%86\%D0\%B8\%D0\%B8}{операциях API} и \href{../07_\%D0\%9F\%D1\%80\%D0\%BE\%D0\%B3\%D1\%80\%D0\%B0\%D0\%BC\%D0\%BC\%D0\%BD\%D1\%8B\%D0\%B5_\%D1\%82\%D1\%80\%D0\%B5\%D0\%B1\%D0\%BE\%D0\%B2\%D0\%B0\%D0\%BD\%D0\%B8\%D1\%8F/DOC-06.2_\%D0\%A1\%D0\%B5\%D1\%81\%D1\%81\%D0\%B8\%D0\%B8_\%D0\%B7\%D0\%B0\%D0\%B4\%D0\%B0\%D1\%87_\%D0\%B8_\%D0\%BF\%D0\%BB\%D0\%B0\%D0\%BD\%D0\%B8\%D1\%80\%D0\%BE\%D0\%B2\%D1\%89\%D0\%B8\%D0\%BA.pdf\#nameddest=\%D0\%BF\%D1\%80\%D0\%B5\%D0\%B4\%D0\%BC\%D0\%B5\%D1\%82\%D0\%BD\%D1\%8B\%D0\%B9-\%D0\%BE\%D1\%82\%D0\%B2\%D0\%B5\%D1\%82}{предметных форматах}. Программная обработка по этому контракту закреплена в \href{../07_\%D0\%9F\%D1\%80\%D0\%BE\%D0\%B3\%D1\%80\%D0\%B0\%D0\%BC\%D0\%BC\%D0\%BD\%D1\%8B\%D0\%B5_\%D1\%82\%D1\%80\%D0\%B5\%D0\%B1\%D0\%BE\%D0\%B2\%D0\%B0\%D0\%BD\%D0\%B8\%D1\%8F/DOC-06.2_\%D0\%A1\%D0\%B5\%D1\%81\%D1\%81\%D0\%B8\%D0\%B8_\%D0\%B7\%D0\%B0\%D0\%B4\%D0\%B0\%D1\%87_\%D0\%B8_\%D0\%BF\%D0\%BB\%D0\%B0\%D0\%BD\%D0\%B8\%D1\%80\%D0\%BE\%D0\%B2\%D1\%89\%D0\%B8\%D0\%BA.pdf\#nameddest=swr-ses-022-\%D1\%84\%D0\%BE\%D1\%80\%D0\%BC\%D0\%B0\%D1\%82\%D1\%8B-\%D0\%B8-\%D0\%BA\%D0\%B0\%D1\%82\%D0\%B5\%D0\%B3\%D0\%BE\%D1\%80\%D0\%B8\%D0\%B8-api}{требовании форматов и категорий API}.

Основание: раздел \texttt{{[}TS.REQ.INT{]}}; область документирования системы.

Проверка: данные, выданные при создании сессии, позволяют однозначно выбрать задачу и сформировать ответ поддерживаемого типа; результат относится к этой задаче. Детальные поля определяет \href{../07_\%D0\%9F\%D1\%80\%D0\%BE\%D0\%B3\%D1\%80\%D0\%B0\%D0\%BC\%D0\%BC\%D0\%BD\%D1\%8B\%D0\%B5_\%D1\%82\%D1\%80\%D0\%B5\%D0\%B1\%D0\%BE\%D0\%B2\%D0\%B0\%D0\%BD\%D0\%B8\%D1\%8F/DOC-06.2_\%D0\%A1\%D0\%B5\%D1\%81\%D1\%81\%D0\%B8\%D0\%B8_\%D0\%B7\%D0\%B0\%D0\%B4\%D0\%B0\%D1\%87_\%D0\%B8_\%D0\%BF\%D0\%BB\%D0\%B0\%D0\%BD\%D0\%B8\%D1\%80\%D0\%BE\%D0\%B2\%D1\%89\%D0\%B8\%D0\%BA.pdf\#nameddest=doc-062-\%D0\%BA\%D0\%BE\%D0\%BD\%D1\%82\%D1\%80\%D0\%B0\%D0\%BA\%D1\%82\%D1\%8B-\%D1\%81\%D0\%B5\%D1\%81\%D1\%81\%D0\%B8\%D0\%B9-\%D0\%B7\%D0\%B0\%D0\%B4\%D0\%B0\%D1\%87}{программный контракт}.

\subsection{SR-SYS-009. Различение отказов}\label{sr-sys-009-ux440ux430ux437ux43bux438ux447ux435ux43dux438ux435-ux43eux442ux43aux430ux437ux43eux432}

\hypertarget{sr-sys-009-различение-отказов}{}

Для состояния задачи имеет значение причина отказа. Неверное решение даёт предметный итог и расходует попытку. Неверно оформленный запрос или технический сбой этого не устанавливает. Просрочка отдельно завершает возможность решения по времени. Если внешняя часть сведёт всё к одному отрицательному ответу, клиент и сервер потеряют различие между этими событиями.

\textbf{Требование.} Система должна отличать неверное предметное решение от отказа в допуске, ошибки формата и технического сбоя проверки.

Клиент определяет категорию по машинному коду, а текст сообщения использует как пояснение. Значение \texttt{isOk\ =\ false} относится только к сохранённому неверному предметному решению. Ошибки допуска и инфраструктуры передаются иначе, по \href{../07_\%D0\%9F\%D1\%80\%D0\%BE\%D0\%B3\%D1\%80\%D0\%B0\%D0\%BC\%D0\%BC\%D0\%BD\%D1\%8B\%D0\%B5_\%D1\%82\%D1\%80\%D0\%B5\%D0\%B1\%D0\%BE\%D0\%B2\%D0\%B0\%D0\%BD\%D0\%B8\%D1\%8F/DOC-06.2_\%D0\%A1\%D0\%B5\%D1\%81\%D1\%81\%D0\%B8\%D0\%B8_\%D0\%B7\%D0\%B0\%D0\%B4\%D0\%B0\%D1\%87_\%D0\%B8_\%D0\%BF\%D0\%BB\%D0\%B0\%D0\%BD\%D0\%B8\%D1\%80\%D0\%BE\%D0\%B2\%D1\%89\%D0\%B8\%D0\%BA.pdf\#nameddest=\%D0\%BA\%D0\%BE\%D0\%B4\%D1\%8B-\%D0\%B8-\%D0\%BF\%D1\%80\%D0\%B5\%D0\%B4\%D1\%81\%D1\%82\%D0\%B0\%D0\%B2\%D0\%BB\%D0\%B5\%D0\%BD\%D0\%B8\%D0\%B5-\%D0\%BE\%D1\%88\%D0\%B8\%D0\%B1\%D0\%BE\%D0\%BA}{таблице ошибок}. Сохранение категорий между компонентами раскрывает \href{../07_\%D0\%9F\%D1\%80\%D0\%BE\%D0\%B3\%D1\%80\%D0\%B0\%D0\%BC\%D0\%BC\%D0\%BD\%D1\%8B\%D0\%B5_\%D1\%82\%D1\%80\%D0\%B5\%D0\%B1\%D0\%BE\%D0\%B2\%D0\%B0\%D0\%BD\%D0\%B8\%D1\%8F/DOC-06_\%D0\%9F\%D1\%80\%D0\%BE\%D0\%B3\%D1\%80\%D0\%B0\%D0\%BC\%D0\%BC\%D0\%BD\%D1\%8B\%D0\%B5_\%D1\%82\%D1\%80\%D0\%B5\%D0\%B1\%D0\%BE\%D0\%B2\%D0\%B0\%D0\%BD\%D0\%B8\%D1\%8F.pdf\#nameddest=swr-sys-007-\%D0\%BA\%D0\%B0\%D1\%82\%D0\%B5\%D0\%B3\%D0\%BE\%D1\%80\%D0\%B8\%D1\%8F-\%D1\%80\%D0\%B5\%D0\%B7\%D1\%83\%D0\%BB\%D1\%8C\%D1\%82\%D0\%B0\%D1\%82\%D0\%B0-\%D0\%BE\%D0\%BF\%D0\%B5\%D1\%80\%D0\%B0\%D1\%86\%D0\%B8\%D0\%B8}{общее программное требование результата операции}.

Основание: правила попытки и ошибок в карточке домена; SR-SES-003, SR-SES-009, SR-SES-010.

Проверка: только зафиксированный предметный вердикт «неверно» расходует попытку как неверное решение. Остальные категории сообщаются отдельно; просрочка учитывается по правилу срока. HTTP-коды и схема ошибок заданы в контракте сессий.

\subsection{Границы обмена между частями}\label{ux433ux440ux430ux43dux438ux446ux44b-ux43eux431ux43cux435ux43dux430-ux43cux435ux436ux434ux443-ux447ux430ux441ux442ux44fux43cux438}

\hypertarget{границы-обмена-между-частями}{}

Связь идентификаторов проверяется на всём пользовательском пути. При создании задачи система устанавливает её принадлежность сессии; при ответе использует эту же принадлежность для проверки владельца. Результат должен относиться к выбранной задаче. Несогласованность хотя бы на одном из этих переходов нарушает общий контракт даже при правильной работе отдельного проверяющего.

У конфигурации и созданной сессии разные обязанности. Конфигурация определяет доступный состав нового набора. Созданная сессия хранит уже выданные задачи и условия. Это разделение позволяет проверять ответ по тому условию, которое получил пользователь, независимо от дальнейшего изменения настроек группы.

Профильная проверка получает условие конкретной задачи и её ответ. Домен сессий отвечает за выбор допустимого обращения и сохранение результата; профильный домен --- за правильность предметного вердикта. Требование к общему контракту не подменяет правила криптографического алгоритма и не назначает допустимые параметры за его владельца.

\subsection{Согласованность при сбоях}\label{ux441ux43eux433ux43bux430ux441ux43eux432ux430ux43dux43dux43eux441ux442ux44c-ux43fux440ux438-ux441ux431ux43eux44fux445}

\hypertarget{согласованность-при-сбоях}{}

Для пользователя подтверждение результата должно означать согласованное изменение задачи и сессии. Успешный вызов предметного проверяющего сам по себе ещё не завершает операцию: результат необходимо сохранить. Если сохранение не подтверждено, система сообщает технический отказ, а не выдаёт неподтверждённый предметный итог как завершённую операцию.

Сохранность завершённых данных и восстановление прерванного вычисления рассматриваются отдельно. Для первого случая определены сохраняемые сведения и неизменность срока после перезапуска. Для второго перед вычислением сохраняется запись принятого ответа. После запуска приложения незавершённые записи продолжают обрабатываться в пределах прежнего срока. Результат фиксируется однократно вместе с состояниями задачи и сессии.

Общая проверка интеграции должна сопоставлять результат операции с данными всех участвующих частей. Для ответа это личность пользователя, принадлежность задачи, сохранённое условие, категория результата и итог сессии. Для контроля времени --- общий срок и факт своевременного поступления ответа. Проверка только одного HTTP-ответа не показывает, сохранена ли согласованность этих сведений.

Редакция 07.11.2026: развёрнуты пояснения каждого требования, добавлены примеры и точные ссылки на связанные разделы. Формулировки обязательств и идентификаторы сохранены.

\section{DOC-05. Требования к ранцевым задачам}\label{doc-05-ux442ux440ux435ux431ux43eux432ux430ux43dux438ux44f-ux43a-ux440ux430ux43dux446ux435ux432ux44bux43c-ux437ux430ux434ux430ux447ux430ux43c}

\hypertarget{doc-05-требования-к-ранцевым-задачам}{}

Раздел \hyperref[doc-05-ux441ux438ux441ux442ux435ux43cux43dux44bux435-ux442ux440ux435ux431ux43eux432ux430ux43dux438ux44f]{DOC-05}, рабочая версия 0.1 от 07.11.2026. Автор: Георгиев Д.Д. Все требования раздела имеют обязательный приоритет. Предметные основания: \href{https://github.com/Kokorev-VD/mephi-samolety/blob/main/notes/domains/backpack.md}{карточка домена} и указанные ниже требования \href{../02_\%D0\%A2\%D0\%B5\%D1\%85\%D0\%BD\%D0\%B8\%D1\%87\%D0\%B5\%D1\%81\%D0\%BA\%D0\%BE\%D0\%B5_\%D0\%B7\%D0\%B0\%D0\%B4\%D0\%B0\%D0\%BD\%D0\%B8\%D0\%B5/\%D0\%A2\%D0\%B5\%D1\%85\%D0\%BD\%D0\%B8\%D1\%87\%D0\%B5\%D1\%81\%D0\%BA\%D0\%BE\%D0\%B5_\%D0\%B7\%D0\%B0\%D0\%B4\%D0\%B0\%D0\%BD\%D0\%B8\%D0\%B5.pdf}{ТЗ}. Состав и формат полей ответа приведены в \href{../07_\%D0\%9F\%D1\%80\%D0\%BE\%D0\%B3\%D1\%80\%D0\%B0\%D0\%BC\%D0\%BC\%D0\%BD\%D1\%8B\%D0\%B5_\%D1\%82\%D1\%80\%D0\%B5\%D0\%B1\%D0\%BE\%D0\%B2\%D0\%B0\%D0\%BD\%D0\%B8\%D1\%8F/DOC-06.2_\%D0\%A1\%D0\%B5\%D1\%81\%D1\%81\%D0\%B8\%D0\%B8_\%D0\%B7\%D0\%B0\%D0\%B4\%D0\%B0\%D1\%87_\%D0\%B8_\%D0\%BF\%D0\%BB\%D0\%B0\%D0\%BD\%D0\%B8\%D1\%80\%D0\%BE\%D0\%B2\%D1\%89\%D0\%B8\%D0\%BA.pdf\#nameddest=\%D1\%80\%D0\%B0\%D0\%BD\%D1\%86\%D0\%B5\%D0\%B2\%D0\%B0\%D1\%8F-\%D0\%B7\%D0\%B0\%D0\%B4\%D0\%B0\%D1\%87\%D0\%B0}{контракте сессий}.

\subsection{SR-KNAP-001. Выбор варианта ранцевой задачи}\label{sr-knap-001-ux432ux44bux431ux43eux440-ux432ux430ux440ux438ux430ux43dux442ux430-ux440ux430ux43dux446ux435ux432ux43eux439-ux437ux430ux434ux430ux447ux438}

\hypertarget{sr-knap-001-выбор-варианта-ранцевой-задачи}{}

Ранцевая криптосистема существует в двух учебных вариантах: на супервозрастающей последовательности и на последовательности степеней. Состав набора задаёт конфигурация группы, но конкретный математический вариант каждого ранцевого задания определяет система при генерации. Пользователь не передаёт серверу свой вариант: он получает условие уже выбранного вида и восстанавливает параметры именно этой криптосистемы. Порядок выбора конфигурации и создания набора определён в \hyperref[sr-ses-001-ux441ux43eux437ux434ux430ux43dux438ux435-ux43dux430ux431ux43eux440ux430-ux43fux43e-ux43aux43eux43dux444ux438ux433ux443ux440ux430ux446ux438ux438-ux433ux440ux443ux43fux43fux44b]{создании сессии}, а роль предметных генераторов --- в \hyperref[sr-ses-004-ux432ux44bux431ux43eux440-ux43fux440ux435ux434ux43cux435ux442ux43dux43eux439-ux43fux440ux43eux432ux435ux440ux43aux438]{выборе предметной проверки}.

\textbf{Требование.} Система должна выбирать вариант ранцевой задачи, выдаваемой в сессии, из поддерживаемых вариантов --- на супервозрастающей последовательности и на степенях --- и формировать условие по схеме выбранного варианта.

Если конфигурация включает ранцевую задачу и задачу Эль-Гамаля, ранцевое условие строится одним из двух поддерживаемых вариантов, а задание Эль-Гамаля --- собственной схемой своего домена. Неподдерживаемый вариант не подменяется другим и не создаёт частичной сессии: это соответствует \href{../07_\%D0\%9F\%D1\%80\%D0\%BE\%D0\%B3\%D1\%80\%D0\%B0\%D0\%BC\%D0\%BC\%D0\%BD\%D1\%8B\%D0\%B5_\%D1\%82\%D1\%80\%D0\%B5\%D0\%B1\%D0\%BE\%D0\%B2\%D0\%B0\%D0\%BD\%D0\%B8\%D1\%8F/DOC-06.2_\%D0\%A1\%D0\%B5\%D1\%81\%D1\%81\%D0\%B8\%D0\%B8_\%D0\%B7\%D0\%B0\%D0\%B4\%D0\%B0\%D1\%87_\%D0\%B8_\%D0\%BF\%D0\%BB\%D0\%B0\%D0\%BD\%D0\%B8\%D1\%80\%D0\%BE\%D0\%B2\%D1\%89\%D0\%B8\%D0\%BA.pdf\#nameddest=swr-ses-005-\%D1\%81\%D0\%BE\%D1\%85\%D1\%80\%D0\%B0\%D0\%BD\%D0\%B5\%D0\%BD\%D0\%B8\%D0\%B5-\%D1\%81\%D0\%BE\%D0\%B7\%D0\%B4\%D0\%B0\%D0\%BD\%D0\%BD\%D0\%BE\%D0\%B3\%D0\%BE-\%D0\%BD\%D0\%B0\%D0\%B1\%D0\%BE\%D1\%80\%D0\%B0}{правилу сохранения полного набора}.

Основание: \texttt{[TS.\allowbreak{}REQ.\allowbreak{}FUNC.\allowbreak{}SESSION.\allowbreak{}CREATE]}, \texttt{[TS.\allowbreak{}REQ.\allowbreak{}FUNC.\allowbreak{}PROB.\allowbreak{}BACKPACK.\allowbreak{}SI]}; сценарий генерации карточки домена.

Проверка: в сессии появляются только задачи поддерживаемых вариантов; схема условия каждой задачи соответствует её варианту; выбор варианта определяется системой и не зависит от параметров пользователя.

\subsection{SR-KNAP-002. Частично скрытое условие}\label{sr-knap-002-ux447ux430ux441ux442ux438ux447ux43dux43e-ux441ux43aux440ux44bux442ux43eux435-ux443ux441ux43bux43eux432ux438ux435}

\hypertarget{sr-knap-002-частично-скрытое-условие}{}

Учебная цель ранцевого задания --- восстановить криптосистему по неполным данным. Поэтому условие не повторяет эталон: часть элементов лёгкого ранца скрыта, а множитель преобразования может отсутствовать вовсе. Студент достраивает недостающее, используя математические свойства последовательности, и затем завершает построение всех величин. Полное эталонное решение с открытыми данными выдаётся отдельно редакционной генерацией по SR-KNAP-008.

\textbf{Требование.} Система должна выдавать условие ранцевой задачи с частично скрытым лёгким ранцем; множитель преобразования может не включаться в условие. Условие должно быть достаточным для восстановления всех параметров криптосистемы.

Восстановление опирается на свойство последовательности: для супервозрастающего варианта каждый элемент строго больше суммы предыдущих, для варианта на степенях элементы являются степенями основания. Скрытие не должно делать задачу неразрешимой: по выданным фрагментам и тяжёлому ранцу восстанавливаются лёгкий ранец, модуль, множитель, обратный множитель, шифртекст и расшифрованное сообщение. Формат выданного условия входит в \href{../07_\%D0\%9F\%D1\%80\%D0\%BE\%D0\%B3\%D1\%80\%D0\%B0\%D0\%BC\%D0\%BC\%D0\%BD\%D1\%8B\%D0\%B5_\%D1\%82\%D1\%80\%D0\%B5\%D0\%B1\%D0\%BE\%D0\%B2\%D0\%B0\%D0\%BD\%D0\%B8\%D1\%8F/DOC-06.2_\%D0\%A1\%D0\%B5\%D1\%81\%D1\%81\%D0\%B8\%D0\%B8_\%D0\%B7\%D0\%B0\%D0\%B4\%D0\%B0\%D1\%87_\%D0\%B8_\%D0\%BF\%D0\%BB\%D0\%B0\%D0\%BD\%D0\%B8\%D1\%80\%D0\%BE\%D0\%B2\%D1\%89\%D0\%B8\%D0\%BA.pdf\#nameddest=\%D0\%BF\%D1\%80\%D0\%B5\%D0\%B4\%D1\%81\%D1\%82\%D0\%B0\%D0\%B2\%D0\%BB\%D0\%B5\%D0\%BD\%D0\%B8\%D0\%B5-\%D1\%81\%D0\%B5\%D1\%81\%D1\%81\%D0\%B8\%D0\%B8-\%D0\%B8-\%D0\%B7\%D0\%B0\%D0\%B4\%D0\%B0\%D1\%87\%D0\%B8}{представление задачи} контракта сессий.

Основание: \texttt{[TS.\allowbreak{}REQ.\allowbreak{}FUNC.\allowbreak{}PROB.\allowbreak{}BACKPACK.\allowbreak{}SI]}; сценарий генерации карточки домена.

Проверка: по выданному условию восстановимы лёгкий ранец, модуль, множитель, тяжёлый ранец, обратный множитель, шифртекст и расшифрованное сообщение.

\subsection{SR-KNAP-003. Состав ответа ранцевой задачи}\label{sr-knap-003-ux441ux43eux441ux442ux430ux432-ux43eux442ux432ux435ux442ux430-ux440ux430ux43dux446ux435ux432ux43eux439-ux437ux430ux434ux430ux447ux438}

\hypertarget{sr-knap-003-состав-ответа-ранцевой-задачи}{}

Ответ на ранцевое задание --- это не отдельное число, а восстановленная криптосистема целиком. Проверяющий сопоставляет каждую переданную величину с ожидаемой, поэтому пропуск поля делает проверку невозможной. Состав полей закреплён \href{../07_\%D0\%9F\%D1\%80\%D0\%BE\%D0\%B3\%D1\%80\%D0\%B0\%D0\%BC\%D0\%BC\%D0\%BD\%D1\%8B\%D0\%B5_\%D1\%82\%D1\%80\%D0\%B5\%D0\%B1\%D0\%BE\%D0\%B2\%D0\%B0\%D0\%BD\%D0\%B8\%D1\%8F/DOC-06.2_\%D0\%A1\%D0\%B5\%D1\%81\%D1\%81\%D0\%B8\%D0\%B8_\%D0\%B7\%D0\%B0\%D0\%B4\%D0\%B0\%D1\%87_\%D0\%B8_\%D0\%BF\%D0\%BB\%D0\%B0\%D0\%BD\%D0\%B8\%D1\%80\%D0\%BE\%D0\%B2\%D1\%89\%D0\%B8\%D0\%BA.pdf\#nameddest=\%D1\%80\%D0\%B0\%D0\%BD\%D1\%86\%D0\%B5\%D0\%B2\%D0\%B0\%D1\%8F-\%D0\%B7\%D0\%B0\%D0\%B4\%D0\%B0\%D1\%87\%D0\%B0}{контрактом сессий}: размерность, вариант, исходное сообщение, лёгкий ранец, множитель, тяжёлый ранец, зашифрованное и расшифрованное сообщения, модуль и обратный множитель. Настоящий раздел подтверждает этот состав как владелец типа.

\textbf{Требование.} Ответ на ранцевую задачу должен содержать восстановленные значения всех величин криптосистемы: размерность, вариант, исходное сообщение, лёгкий ранец, множитель, тяжёлый ранец, зашифрованное и расшифрованное сообщения, модуль и обратный множитель. Указанный в ответе вариант не должен подменять вариант выбранной задачи.

Поля передаются в установленном контрактом виде: размерность --- положительным целым числом, сообщения --- списками логических значений, целые величины произвольной длины --- десятичными строками, ранцы --- списками таких строк. Ответ с полным набором полей допускается к предметной проверке; категория ошибки формы определяется \hyperref[sr-ses-009-ux43eux448ux438ux431ux43aux430-ux444ux43eux440ux43cux430ux442ux430]{правилом ошибки формата} домена сессий. Значение поля варианта сверяется с вариантом сохранённого условия и не переключает проверку на другой тип.

Основание: \texttt{[TS.\allowbreak{}REQ.\allowbreak{}FUNC.\allowbreak{}SESSION.\allowbreak{}SUBMIT]}, \texttt{[TS.\allowbreak{}REQ.\allowbreak{}FUNC.\allowbreak{}PROB.\allowbreak{}BACKPACK.\allowbreak{}SI]}; поля предметного ответа в контракте сессий.

Проверка: ответ с полным набором величин допускается к предметной проверке; значение поля варианта не переключает проверку на другой вариант. Категории отказа до предметной проверки определяются SR-SES-009.

\subsection{SR-KNAP-004. Предметная проверка решения}\label{sr-knap-004-ux43fux440ux435ux434ux43cux435ux442ux43dux430ux44f-ux43fux440ux43eux432ux435ux440ux43aux430-ux440ux435ux448ux435ux43dux438ux44f}

\hypertarget{sr-knap-004-предметная-проверка-решения}{}

Предметная проверка ранцевого задания проходит по всем восстановленным величинам, а не по итоговому сообщению. Каждая величина вычисляется из условия и сравнивается с переданной. Перечень охватывает длину и элементы сообщений и ранцев, сверхвозрастаемость лёгкого ранца, соотношение модуля с суммой лёгкого ранца, взаимную простоту множителя с модулем, тяжёлый ранец, обратный множитель, шифртекст и расшифрованное сообщение. Такой состав следует из \href{../02_\%D0\%A2\%D0\%B5\%D1\%85\%D0\%BD\%D0\%B8\%D1\%87\%D0\%B5\%D1\%81\%D0\%BA\%D0\%BE\%D0\%B5_\%D0\%B7\%D0\%B0\%D0\%B4\%D0\%B0\%D0\%BD\%D0\%B8\%D0\%B5/\%D0\%A2\%D0\%B5\%D1\%85\%D0\%BD\%D0\%B8\%D1\%87\%D0\%B5\%D1\%81\%D0\%BA\%D0\%BE\%D0\%B5_\%D0\%B7\%D0\%B0\%D0\%B4\%D0\%B0\%D0\%BD\%D0\%B8\%D0\%B5.pdf}{требования ТЗ к проверке восстанавливаемых величин} и сценария проверки карточки домена.

\textbf{Требование.} Система должна проверять решение ранцевой задачи по всем восстановимым величинам: длине и элементам сообщения и ранцев, сверхвозрастаемости лёгкого ранца, соотношению модуля с суммой лёгкого ранца, взаимной простоте множителя с модулем, тяжёлому ранцу, обратному множителю, шифртексту и расшифрованному сообщению.

Расшифрованное сообщение проверяется расшифровкой жадным алгоритмом по супервозрастающей последовательности: элементы выбираются от старших к младшим, и результат сравнивается с переданным сообщением. Расхождение любой величины с ожидаемой приводит к неверному вердикту; полностью корректное решение --- к верному. Результат операции и состояние задачи согласуются по \hyperref[sr-ses-005-ux440ux435ux437ux443ux43bux44cux442ux430ux442-ux437ux430ux434ux430ux447ux438]{результату задачи}, а длительность проверки ограничена \hyperref[sr-ses-012-ux43eux433ux440ux430ux43dux438ux447ux435ux43dux43dux430ux44f-ux434ux43bux438ux442ux435ux43bux44cux43dux43eux441ux442ux44c-ux43eux431ux440ux430ux431ux43eux442ux43aux438]{пределом обработки} домена сессий.

Основание: \texttt{[TS.\allowbreak{}REQ.\allowbreak{}FUNC.\allowbreak{}PROB.\allowbreak{}BACKPACK.\allowbreak{}SI]}; сценарий проверки карточки домена.

Проверка: расхождение любой величины с ожидаемой приводит к неверному вердикту; полностью корректное решение --- к верному.

\subsection{SR-KNAP-005. Диагностика неверного решения}\label{sr-knap-005-ux434ux438ux430ux433ux43dux43eux441ux442ux438ux43aux430-ux43dux435ux432ux435ux440ux43dux43eux433ux43e-ux440ux435ux448ux435ux43dux438ux44f}

\hypertarget{sr-knap-005-диагностика-неверного-решения}{}

Неверный ответ без пояснения не даёт студенту материала для работы над ошибкой. Предметная проверка сообщает, какая величина не прошла сравнение первой. Это пояснение относится к завершённой неверной попытке и не заменяет категорию отказа: транспортные и технические ошибки передаются отдельными кодами по \href{../07_\%D0\%9F\%D1\%80\%D0\%BE\%D0\%B3\%D1\%80\%D0\%B0\%D0\%BC\%D0\%BC\%D0\%BD\%D1\%8B\%D0\%B5_\%D1\%82\%D1\%80\%D0\%B5\%D0\%B1\%D0\%BE\%D0\%B2\%D0\%B0\%D0\%BD\%D0\%B8\%D1\%8F/DOC-06.2_\%D0\%A1\%D0\%B5\%D1\%81\%D1\%81\%D0\%B8\%D0\%B8_\%D0\%B7\%D0\%B0\%D0\%B4\%D0\%B0\%D1\%87_\%D0\%B8_\%D0\%BF\%D0\%BB\%D0\%B0\%D0\%BD\%D0\%B8\%D1\%80\%D0\%BE\%D0\%B2\%D1\%89\%D0\%B8\%D0\%BA.pdf\#nameddest=\%D0\%BA\%D0\%BE\%D0\%B4\%D1\%8B-\%D0\%B8-\%D0\%BF\%D1\%80\%D0\%B5\%D0\%B4\%D1\%81\%D1\%82\%D0\%B0\%D0\%B2\%D0\%BB\%D0\%B5\%D0\%BD\%D0\%B8\%D0\%B5-\%D0\%BE\%D1\%88\%D0\%B8\%D0\%B1\%D0\%BE\%D0\%BA}{контракту ошибок}.

\textbf{Требование.} При неверном решении система должна сообщать, какая величина не прошла проверку первой.

Диагностика указывает конкретную величину, например неверный тяжёлый ранец или нарушенную взаимную простоту множителя. Верный вердикт диагностикой не сопровождается: поле пояснения в положительном результате отсутствует. Текст пояснения не влияет на состояние задачи и попытку --- они определяются требованиями домена сессий, а не формулировкой сообщения. Это применение общего \hyperref[sr-sys-009-ux440ux430ux437ux43bux438ux447ux435ux43dux438ux435-ux43eux442ux43aux430ux437ux43eux432]{различения отказов}.

Основание: \texttt{[TS.\allowbreak{}REQ.\allowbreak{}FUNC.\allowbreak{}SESSION.\allowbreak{}SUBMIT]} (диагностическое сообщение при ошибке).

Проверка: сообщение указывает конкретную величину; верный вердикт не сопровождается диагностикой ошибки. Сообщение не заменяет категорию результата по SR-SYS-009.

\subsection{SR-KNAP-006. Вариант на степенях}\label{sr-knap-006-ux432ux430ux440ux438ux430ux43dux442-ux43dux430-ux441ux442ux435ux43fux435ux43dux44fux445}

\hypertarget{sr-knap-006-вариант-на-степенях}{}

Второй вариант ранцевой криптосистемы отличается построением лёгкого ранца: его элементы являются степенями выбранного основания. Такой ранец автоматически обладает сверхвозрастаемостью, что и позволяет применять общую схему проверки. Студент, решающий задание на степенях, сначала распознаёт структуру последовательности, затем восстанавливает остальные величины по той же цепочке, что и в супервозрастающем варианте.

\textbf{Требование.} Для варианта «на степенях» система должна строить лёгкий ранец из степеней выбранного основания и проверять решение по той же схеме величин, что и для супервозрастающего варианта.

Состав проверяемых величин совпадает со схемой SR-KNAP-004. Фактическое отсутствие проверки расшифрованного сообщения для этого варианта в текущей реализации отмечено в DOC-06.3 как известная особенность и вынесено на решение к следующей версии документа.

Основание: \texttt{[TS.\allowbreak{}REQ.\allowbreak{}FUNC.\allowbreak{}PROB.\allowbreak{}BACKPACK.\allowbreak{}DEG]}.

Проверка: элементы лёгкого ранца соответствуют степеням основания; состав проверок совпадает со схемой SR-KNAP-004.

\subsection{SR-KNAP-007. Повторный ответ после итогового состояния}\label{sr-knap-007-ux43fux43eux432ux442ux43eux440ux43dux44bux439-ux43eux442ux432ux435ux442-ux43fux43eux441ux43bux435-ux438ux442ux43eux433ux43eux432ux43eux433ux43e-ux441ux43eux441ux442ux43eux44fux43dux438ux44f}

\hypertarget{sr-knap-007-повторный-ответ-после-итогового-состояния}{}

После зафиксированного вердикта повторная отправка решения не создаёт новой попытки. Это правило домена сессий сформулировано в \hyperref[sr-ses-008-ux43eux434ux43dux430-ux43fux440ux435ux434ux43cux435ux442ux43dux430ux44f-ux43fux43eux43fux44bux442ux43aux430]{требовании единственной попытки}; ранцевой домен применяет его к своему обработчику. Повтор может произойти из-за повторного нажатия, повторной отправки сетевого клиента или нового соединения --- во всех случаях прежний результат сохраняется.

\textbf{Требование.} После итогового состояния задачи система должна отклонять повторный ответ ранцевой задачи без предметной проверки, сохраняя прежний результат.

Отклонение происходит до математического вычисления: проверяющий не запускается, а задача и сессия не изменяются. Внешняя категория отказа определена контрактом сессий (\texttt{TASK\_FINALIZED}); для восстановления отображения студент использует \hyperref[sr-ses-014-ux43fux43eux43bux443ux447ux435ux43dux438ux435-ux441ux43eux445ux440ux430ux43dux451ux43dux43dux43eux433ux43e-ux440ux435ux437ux443ux43bux44cux442ux430ux442ux430]{чтение сохранённого результата}, а не повторную отправку.

Основание: \texttt{[TS.\allowbreak{}REQ.\allowbreak{}FUNC.\allowbreak{}SESSION.\allowbreak{}SUBMIT]}; правило попытки карточки домена.

Проверка: повторная отправка не создаёт нового предметного вердикта и не изменяет состояние задачи; согласуется с SR-SES-008.

\subsection{SR-KNAP-008. Редакционная задача по варианту}\label{sr-knap-008-ux440ux435ux434ux430ux43aux446ux438ux43eux43dux43dux430ux44f-ux437ux430ux434ux430ux447ux430-ux43fux43e-ux432ux430ux440ux438ux430ux43dux442ux443}

\hypertarget{sr-knap-008-редакционная-задача-по-варианту}{}

Редакционная генерация предназначена для самоподготовки и подготовки методических материалов. В отличие от задания в сессии, она выдаёт эталон с полным решением: открытым лёгким ранцем, множителем, модулем, тяжёлым ранцем, обратным множителем и шифртекстом. По эталону студент прослеживает всю цепочку построения до расшифрованного сообщения.

\textbf{Требование.} Система должна выдавать редакционную (эталонную) ранцевую задачу выбранного варианта с полным решением --- всеми величинами криптосистемы от лёгкого ранца до расшифрованного сообщения.

Вариант задаётся запросом; остальные параметры, не указанные явно, выбираются системой. Редакционные генерации не связаны с сессиями и не влияют на попытки и сроки --- их единственная обязанность перед остальной системой описана в SR-KNAP-010.

Основание: \texttt{[TS.\allowbreak{}REQ.\allowbreak{}FUNC.\allowbreak{}EDITORIAL]}; учебный сценарий карточки домена.

Проверка: эталон содержит полный набор величин; вариант соответствует указанному в запросе.

\subsection{SR-KNAP-009. Редакционная задача с параметрами}\label{sr-knap-009-ux440ux435ux434ux430ux43aux446ux438ux43eux43dux43dux430ux44f-ux437ux430ux434ux430ux447ux430-ux441-ux43fux430ux440ux430ux43cux435ux442ux440ux430ux43cux438}

\hypertarget{sr-knap-009-редакционная-задача-с-параметрами}{}

Преподаватель может задать собственные параметры примера: размерность, вариант, сообщение, лёгкий ранец и множитель. Система строит эталон по этим данным, проверяя математическую допустимость. Взаимная простота множителя с модулем --- необходимое условие обратного преобразования, поэтому её нарушение отклоняется до построения.

\textbf{Требование.} Система должна принимать параметры редакционной генерации --- размерность, вариант, сообщение, лёгкий ранец, множитель --- и отклонять запрос с множителем, не взаимно простым с модулем, сообщая причину.

При отказе задача не выдаётся: клиент получает диагностическое сообщение о причине. Отказ редакционной генерации не связан с сессиями и не расходует никаких попыток.

Основание: \texttt{[TS.\allowbreak{}REQ.\allowbreak{}FUNC.\allowbreak{}EDITORIAL]}; сценарий с параметрами карточки домена.

Проверка: генерация с корректными параметрами завершается эталоном; при невзаимно простом множителе --- отказ с диагностикой и без задачи.

\subsection{SR-KNAP-010. Редакционные генерации без сохранения}\label{sr-knap-010-ux440ux435ux434ux430ux43aux446ux438ux43eux43dux43dux44bux435-ux433ux435ux43dux435ux440ux430ux446ux438ux438-ux431ux435ux437-ux441ux43eux445ux440ux430ux43dux435ux43dux438ux44f}

\hypertarget{sr-knap-010-редакционные-генерации-без-сохранения}{}

Редакционные примеры не являются учебными попытками и не должны попадать в данные сессий. Каждый запрос создаёт новый пример на лету; сохранение сделало бы повторный запрос выдачей прежнего эталона и связало бы самоподготовку с хранилищем задач.

\textbf{Требование.} Система не должна записывать результаты редакционных генераций в базу данных.

Правило относится к результатам редакционных генераций и не затрагивает сохранение задач, созданных в сессиях по SR-KNAP-001.

Основание: \texttt{[TS.\allowbreak{}REQ.\allowbreak{}FUNC.\allowbreak{}EDITORIAL]}.

Проверка: повторный запрос даёт новый пример; сохранённые данные сессий не пополняются редакционными задачами.

\subsection{Пояснения к вариантам и попытке}\label{ux43fux43eux44fux441ux43dux435ux43dux438ux44f-ux43a-ux432ux430ux440ux438ux430ux43dux442ux430ux43c-ux438-ux43fux43eux43fux44bux442ux43aux435}

\hypertarget{пояснения-к-вариантам-и-попытке}{}

Вариант задания и вариант в ответе студента имеют разное назначение. Первый выбирается системой при генерации и сохраняется вместе с условием; второй передаётся в теле ответа и сверяется с первым. Совпадение необходимо: ответ по схеме другого варианта не может проверить сохранённое условие.

Предметная попытка расходуется только завершённой проверкой с вердиктом. Неверная величина в ответе даёт отрицательный вердикт и завершает попытку; ошибка формы или технический сбой проверки обрабатываются правилами домена сессий без участия ранцевых алгоритмов. Диагностическое сообщение поясняет неверное решение, но не создаёт и не отменяет попытку.

Длительность предметной проверки входит в предел обработки принятого ответа, установленный доменом сессий. Ранцевая проверка по величинам конечна и не требует собственных таймеров: превышение предела обработки обработывается общим правилом технического завершения.

\subsection{Сочетания событий для проверки раздела}\label{ux441ux43eux447ux435ux442ux430ux43dux438ux44f-ux441ux43eux431ux44bux442ux438ux439-ux434ux43bux44f-ux43fux440ux43eux432ux435ux440ux43aux438-ux440ux430ux437ux434ux435ux43bux430}

\hypertarget{сочетания-событий-для-проверки-раздела}{}

Приведённые сочетания раскрывают требования, а не вводят отдельные TC. При подготовке DOC-10 для них нужно задать исходное состояние, действие и наблюдаемый результат.

{\def\LTcaptype{none} % do not increment counter
\begin{longtable}[]{@{}
  >{\raggedright\arraybackslash}p{(\linewidth - 2\tabcolsep) * \real{0.3200}}
  >{\raggedright\arraybackslash}p{(\linewidth - 2\tabcolsep) * \real{0.6800}}@{}}
\toprule\noalign{}
\begin{minipage}[b]{\linewidth}\raggedright
Последовательность событий
\end{minipage} & \begin{minipage}[b]{\linewidth}\raggedright
Ожидаемый результат
\end{minipage} \\
\midrule\noalign{}
\endhead
\bottomrule\noalign{}
\endlastfoot
Создан набор с ранцевой задачей; условие содержит скрытые элементы ранца & Условие восстанавливается до полного набора величин; вариант соответствует схеме \\
Ответ с полным набором величин, одна величина неверна & Отрицательный вердикт; диагностика указывает первую несовпавшую величину \\
Ответ с полным набором величин, все верны & Положительный вердикт; попытка завершена \\
Поле варианта в ответе указывает другой вариант & Ответ не допускается к проверке по чужой схеме; вариант сверяется с условием \\
Повторная отправка после зафиксированного вердикта & Отказ без предметной проверки; прежний результат сохранён \\
Редакционная генерация выполнена дважды с одними параметрами & Каждый запрос даёт новый пример; хранилище сессий не пополняется \\
Редакционная генерация с невзаимно простым множителем & Отказ с диагностикой; пример не выдаётся \\
\end{longtable}
}

История: 07.11.2026 --- подготовлена первая редакция раздела: варианты задач, условие, состав ответа, предметная проверка, редакционные генерации.


\end{document}
